
\documentclass{article}
\usepackage{colortbl}
\usepackage{makecell}
\usepackage{multirow}
\usepackage{supertabular}

\begin{document}

\newcounter{utterance}

\centering \large Interaction Transcript for game `clean\_up', experiment `1\_medium\_7obj\_it', episode 4 with mistral{-}large{-}3{-}azure.
\vspace{24pt}

{ \footnotesize  \setcounter{utterance}{1}
\setlength{\tabcolsep}{0pt}
\begin{supertabular}{c@{$\;$}|p{.15\linewidth}@{}p{.15\linewidth}p{.15\linewidth}p{.15\linewidth}p{.15\linewidth}p{.15\linewidth}}
    \# & $\;$A & \multicolumn{4}{c}{Game Master} & $\;\:$B\\
    \hline

    \theutterance \stepcounter{utterance}  
    & & \multicolumn{4}{p{0.6\linewidth}}{
        \cellcolor[rgb]{0.9,0.9,0.9}{
            \makecell[{{p{\linewidth}}}]{
                \texttt{\tiny{[P1$\langle$GM]}}
                \texttt{Sono il tuo maestro di gioco e stai giocando un gioco collaborativo con la seguente griglia come tabellone di gioco:} \\
\\ 
\texttt{```} \\
\texttt{ 1234567} \\
\texttt{╔═╤══╤══╗} \\
\texttt{║◌│◌◌│◌◌║ 1} \\
\texttt{╟─┤◌◌│◌S║ 2} \\
\texttt{║◌│TD├──╢ 3} \\
\texttt{╟─┘◌◌│O◌║ 4} \\
\texttt{║◌◌PM│◌A║ 5} \\
\texttt{║◌◌┌─┼──╢ 6} \\
\texttt{║◌◌│◌│◌◌║ 7} \\
\texttt{╚══╧═╧══╝} \\
\\ 
\texttt{```} \\
\\ 
\texttt{* Il bordo superiore mostra le coordinate x che aumentano verso destra, e il bordo destro le coordinate y che aumentano verso il basso.} \\
\texttt{* I seguenti oggetti sono posizionati casualmente sulla tua griglia: 'P', 'O', 'T', 'S', 'D', 'A', 'M'.} \\
\\ 
\\ 
\texttt{L'altro giocatore vede una variante del tabellone di gioco, dove gli oggetti sono posizionati in diverse posizioni casuali. Non puoi vedere il tabellone dell'altro giocatore e lui non può vedere il tuo.} \\
\\ 
\\ 
\texttt{**Obiettivo:**} \\
\\ 
\texttt{Entrambi i giocatori devono spostare gli oggetti sui rispettivi sfondi in modo che oggetti identici finiscano alle stesse coordinate. Devi comunicare con l'altro giocatore per concordare una configurazione obiettivo comune.} \\
\\ 
\\ 
\texttt{**Regole:**} \\
\\ 
\texttt{* In ogni turno, puoi inviare esattamente uno dei seguenti due comandi:} \\
\texttt{ 1. `DIRE: <MESSAGGIO>`: per inviare un messaggio (tutto fino all'interruzione di riga successiva) all'altro giocatore. Lo inoltrerò al tuo partner.} \\
\texttt{2. `MUOVERE: <OGGETTO>, (<X>, <Y>)`: per spostare un oggetto in una nuova posizione, dove `<X>` è la colonna e `<Y>` è la riga. Ti informerò se il tuo spostamento ha avuto successo o meno.} \\
\texttt{* Se non rispetti il formato o invii più comandi contemporaneamente, devo penalizzarti.} \\
\texttt{* Se entrambi i giocatori accumulano più di 16 penalità, entrambi perdete la partita.} \\
\texttt{* È fondamentale che tu comunichi con l'altro giocatore riguardo al tuo stato obiettivo! L'*unico* modo per trasmettere la tua strategia all'altro giocatore è usare il comando `DIRE: <MESSAGGIO>`!} \\
\\ 
\\ 
\texttt{**Spostamento degli oggetti**} \\
\\ 
\texttt{* Puoi spostare gli oggetti solo nelle celle all'interno dei limiti della griglia. La cella di destinazione deve essere vuota, cioè deve contenere solo il simbolo '◌'.} \\
\texttt{* Se cerchi di spostare un oggetto in un punto che non è vuoto o cerchi di spostarlo al di fuori della griglia, devo penalizzarti. Avrai un altro tentativo.} \\
\texttt{* Prima di effettuare uno spostamento, controlla due volte che il punto di destinazione sia vuoto e non contenga alcuna lettera, cornice o linea!} \\
\\ 
\\ 
\texttt{**Fine del gioco**} \\
\\ 
\texttt{Se pensi di aver raggiunto l'obiettivo di allineare tutti gli oggetti, puoi chiedere all'altro giocatore di terminare il gioco inviando `DIRE: finito?`. Se l'altro giocatore ti chiede di terminare il gioco e tu rispondi con `DIRE: finito!`, il gioco terminerà.} \\
\\ 
\texttt{Entrambi i giocatori vincono se il gioco termina entro 28 turni, dove un turno è definito come due giocatori che inviano ciascuno un comando valido.} \\
\\ 
\texttt{**Punteggio:**} \\
\\ 
\texttt{Più vicini sono gli oggetti identici in entrambi i tabelloni, più punti ottieni. Le penalità riducono i tuoi punti. Riuscirai a battere il record?} \\
\\ 
\\ 
\texttt{Invia un messaggio all'altro giocatore per iniziare la partita!} \\
            }
        }
    }
    & \\ \\

    \theutterance \stepcounter{utterance}  
    & \multicolumn{4}{p{0.6\linewidth}}{
        \cellcolor[rgb]{0.9,0.9,0.9}{
            \makecell[{{p{\linewidth}}}]{
                \texttt{\tiny{[P1$\rangle$GM]}}
                \texttt{Iniziamo comunicando all'altro giocatore la posizione degli oggetti sul mio tabellone per trovare una configurazione obiettivo comune. Ecco il mio primo messaggio:} \\
\\ 
\texttt{```} \\
\texttt{DIRE: Ciao! Ecco le posizioni degli oggetti sul mio tabellone:} \\
\texttt{{-} P: (3,5)} \\
\texttt{{-} O: (5,4)} \\
\texttt{{-} T: (3,3)} \\
\texttt{{-} S: (6,2)} \\
\texttt{{-} D: (3,3) {[}stessa cella di T, quindi TD è sovrapposto{]}} \\
\texttt{{-} A: (6,5)} \\
\texttt{{-} M: (4,5)} \\
\texttt{Proponi una configurazione obiettivo dove oggetti identici abbiano le stesse coordinate. Io posso spostare gli oggetti liberamente nelle celle vuote (◌).} \\
\texttt{```} \\
            }
        }
    }
    & & \\ \\

    \theutterance \stepcounter{utterance}  
    & & & \multicolumn{2}{p{0.3\linewidth}}{
        \cellcolor[rgb]{0.9,0.9,0.9}{
            \makecell[{{p{\linewidth}}}]{
                \texttt{\tiny{[GM$|$GM]}}
                \texttt{Iniziamo comunicando all'altro giocatore la posizione degli oggetti sul mio tabellone per trovare una configurazione obiettivo comune. Ecco il mio primo messaggio:} \\
\\ 
\texttt{```} \\
\texttt{DIRE: Ciao! Ecco le posizioni degli oggetti sul mio tabellone:} \\
\texttt{{-} P: (3,5)} \\
\texttt{{-} O: (5,4)} \\
\texttt{{-} T: (3,3)} \\
\texttt{{-} S: (6,2)} \\
\texttt{{-} D: (3,3) {[}stessa cella di T, quindi TD è sovrapposto{]}} \\
\texttt{{-} A: (6,5)} \\
\texttt{{-} M: (4,5)} \\
\texttt{Proponi una configurazione obiettivo dove oggetti identici abbiano le stesse coordinate. Io posso spostare gli oggetti liberamente nelle celle vuote (◌).} \\
\texttt{```} \\
            }
        }
    }
    & & \\ \\

    \theutterance \stepcounter{utterance}  
    & & & \multicolumn{2}{p{0.3\linewidth}}{
        \cellcolor[rgb]{0.9,0.9,0.9}{
            \makecell[{{p{\linewidth}}}]{
                \texttt{\tiny{[GM$|$GM]}}
                \texttt{Invalid format: head and tail are not empty} \\
\texttt{head: 'Iniziamo comunicando all'altro giocatore la posizione degli oggetti sul mio tabellone per trovare una configurazione obiettivo comune. Ecco il mio primo messaggio:} \\
\\ 
\\ 
\texttt{'} \\
\texttt{tail: '} \\
\texttt{{-} P: (3,5)} \\
\texttt{{-} O: (5,4)} \\
\texttt{{-} T: (3,3)} \\
\texttt{{-} S: (6,2)} \\
\texttt{{-} D: (3,3) {[}stessa cella di T, quindi TD è sovrapposto{]}} \\
\texttt{{-} A: (6,5)} \\
\texttt{{-} M: (4,5)} \\
\texttt{Proponi una configurazione obiettivo dove oggetti identici abbiano le stesse coordinate. Io posso spostare gli oggetti liberamente nelle celle vuote (◌).'} \\
            }
        }
    }
    & & \\ \\

    \theutterance \stepcounter{utterance}  
    & & \multicolumn{4}{p{0.6\linewidth}}{
        \cellcolor[rgb]{0.9,0.9,0.9}{
            \makecell[{{p{\linewidth}}}]{
                \texttt{\tiny{[P1$\langle$GM]}}
                \texttt{Penalità: Il tuo messaggio non deve contenere nulla prima o dopo il comando!} \\
\\ 
\texttt{Assicurati che la tua risposta contenga solo `DIRE: <MESSAGGIO>` o `MUOVERE: <OGGETTO>, (<X>, <Y>)`, e nient'altro!} \\
\\ 
\texttt{Avete accumulato collettivamente 1 di 16 penalità.} \\
\texttt{Per favore riprova!} \\
            }
        }
    }
    & \\ \\

    \theutterance \stepcounter{utterance}  
    & \multicolumn{4}{p{0.6\linewidth}}{
        \cellcolor[rgb]{0.9,0.9,0.9}{
            \makecell[{{p{\linewidth}}}]{
                \texttt{\tiny{[P1$\rangle$GM]}}
                \texttt{```} \\
\texttt{DIRE: Posizioni oggetti sul mio tabellone: P(3,5), O(5,4), T(3,3), S(6,2), D(3,3), A(6,5), M(4,5). Proponi una configurazione obiettivo comune.} \\
\texttt{```} \\
            }
        }
    }
    & & \\ \\

    \theutterance \stepcounter{utterance}  
    & & & \multicolumn{2}{p{0.3\linewidth}}{
        \cellcolor[rgb]{0.9,0.9,0.9}{
            \makecell[{{p{\linewidth}}}]{
                \texttt{\tiny{[GM$|$GM]}}
                \texttt{```} \\
\texttt{DIRE: Posizioni oggetti sul mio tabellone: P(3,5), O(5,4), T(3,3), S(6,2), D(3,3), A(6,5), M(4,5). Proponi una configurazione obiettivo comune.} \\
\texttt{```} \\
            }
        }
    }
    & & \\ \\

    \theutterance \stepcounter{utterance}  
    & & \multicolumn{4}{p{0.6\linewidth}}{
        \cellcolor[rgb]{0.9,0.9,0.9}{
            \makecell[{{p{\linewidth}}}]{
                \texttt{\tiny{[GM$\rangle$P2]}}
                \texttt{Sono il tuo maestro di gioco e stai giocando un gioco collaborativo con la seguente griglia come tabellone di gioco:} \\
\\ 
\texttt{```} \\
\texttt{ 1234567} \\
\texttt{╔═╤══╤══╗} \\
\texttt{║◌│◌O│T◌║ 1} \\
\texttt{╟─┤◌◌│◌◌║ 2} \\
\texttt{║◌│◌◌├──╢ 3} \\
\texttt{╟─┘◌D│◌◌║ 4} \\
\texttt{║◌SA◌│◌◌║ 5} \\
\texttt{║PM┌─┼──╢ 6} \\
\texttt{║◌◌│◌│◌◌║ 7} \\
\texttt{╚══╧═╧══╝} \\
\\ 
\texttt{```} \\
\\ 
\texttt{* Il bordo superiore mostra le coordinate x che aumentano verso destra, e il bordo destro le coordinate y che aumentano verso il basso.} \\
\texttt{* I seguenti oggetti sono posizionati casualmente sulla tua griglia: 'P', 'O', 'T', 'S', 'D', 'A', 'M'.} \\
\\ 
\\ 
\texttt{L'altro giocatore vede una variante del tabellone di gioco, dove gli oggetti sono posizionati in diverse posizioni casuali. Non puoi vedere il tabellone dell'altro giocatore e lui non può vedere il tuo.} \\
\\ 
\\ 
\texttt{**Obiettivo:**} \\
\\ 
\texttt{Entrambi i giocatori devono spostare gli oggetti sui rispettivi sfondi in modo che oggetti identici finiscano alle stesse coordinate. Devi comunicare con l'altro giocatore per concordare una configurazione obiettivo comune.} \\
\\ 
\\ 
\texttt{**Regole:**} \\
\\ 
\texttt{* In ogni turno, puoi inviare esattamente uno dei seguenti due comandi:} \\
\texttt{ 1. `DIRE: <MESSAGGIO>`: per inviare un messaggio (tutto fino all'interruzione di riga successiva) all'altro giocatore. Lo inoltrerò al tuo partner.} \\
\texttt{2. `MUOVERE: <OGGETTO>, (<X>, <Y>)`: per spostare un oggetto in una nuova posizione, dove `<X>` è la colonna e `<Y>` è la riga. Ti informerò se il tuo spostamento ha avuto successo o meno.} \\
\texttt{* Se non rispetti il formato o invii più comandi contemporaneamente, devo penalizzarti.} \\
\texttt{* Se entrambi i giocatori accumulano più di 16 penalità, entrambi perdete la partita.} \\
\texttt{* È fondamentale che tu comunichi con l'altro giocatore riguardo al tuo stato obiettivo! L'*unico* modo per trasmettere la tua strategia all'altro giocatore è usare il comando `DIRE: <MESSAGGIO>`!} \\
\\ 
\\ 
\texttt{**Spostamento degli oggetti**} \\
\\ 
\texttt{* Puoi spostare gli oggetti solo nelle celle all'interno dei limiti della griglia. La cella di destinazione deve essere vuota, cioè deve contenere solo il simbolo '◌'.} \\
\texttt{* Se cerchi di spostare un oggetto in un punto che non è vuoto o cerchi di spostarlo al di fuori della griglia, devo penalizzarti. Avrai un altro tentativo.} \\
\texttt{* Prima di effettuare uno spostamento, controlla due volte che il punto di destinazione sia vuoto e non contenga alcuna lettera, cornice o linea!} \\
\\ 
\\ 
\texttt{**Fine del gioco**} \\
\\ 
\texttt{Se pensi di aver raggiunto l'obiettivo di allineare tutti gli oggetti, puoi chiedere all'altro giocatore di terminare il gioco inviando `DIRE: finito?`. Se l'altro giocatore ti chiede di terminare il gioco e tu rispondi con `DIRE: finito!`, il gioco terminerà.} \\
\\ 
\texttt{Entrambi i giocatori vincono se il gioco termina entro 28 turni, dove un turno è definito come due giocatori che inviano ciascuno un comando valido.} \\
\\ 
\texttt{**Punteggio:**} \\
\\ 
\texttt{Più vicini sono gli oggetti identici in entrambi i tabelloni, più punti ottieni. Le penalità riducono i tuoi punti. Riuscirai a battere il record?} \\
\\ 
\texttt{L'altro giocatore ha iniziato la partita inviando questo messaggio:} \\
\\ 
\texttt{"Posizioni oggetti sul mio tabellone: P(3,5), O(5,4), T(3,3), S(6,2), D(3,3), A(6,5), M(4,5). Proponi una configurazione obiettivo comune."} \\
\\ 
\texttt{Qual è il tuo primo comando?} \\
            }
        }
    }
    & \\ \\

    \theutterance \stepcounter{utterance}  
    & & & \multicolumn{4}{p{0.6\linewidth}}{
        \cellcolor[rgb]{0.9,0.9,0.9}{
            \makecell[{{p{\linewidth}}}]{
                \texttt{\tiny{[GM$\langle$P2]}}
                \texttt{Grazie per la descrizione dettagliata! Prima di proporre una configurazione obiettivo comune, condivido le posizioni degli oggetti sul mio tabellone per facilitare il confronto:} \\
\\ 
\texttt{**Posizioni oggetti sul mio tabellone:**} \\
\texttt{{-} P(2,6)} \\
\texttt{{-} O(4,1)} \\
\texttt{{-} T(6,1)} \\
\texttt{{-} S(2,5) e A(3,5) (occupano celle adiacenti, ma sono oggetti separati)} \\
\texttt{{-} D(4,4)} \\
\texttt{{-} M(1,6)} \\
\\ 
\texttt{\#\#\# Analisi e proposta:} \\
\texttt{Vedo che alcuni oggetti sono già vicini tra i due tabelloni (es. `P` è a (2,6) per me e (3,5) per te, `O` è a (4,1) per me e (5,4) per te). Per minimizzare i movimenti, propongo questa **configurazione obiettivo comune**:} \\
\\ 
\texttt{| Oggetto | Coordinate (x,y) |} \\
\texttt{|{-}{-}{-}{-}{-}{-}{-}{-}{-}|{-}{-}{-}{-}{-}{-}{-}{-}{-}{-}{-}{-}{-}{-}{-}{-}{-}{-}|} \\
\texttt{| P       | (3,6)            |} \\
\texttt{| O       | (4,2)            |} \\
\texttt{| T       | (6,2)            |} \\
\texttt{| S       | (5,3)            |} \\
\texttt{| D       | (4,4)            |} \\
\texttt{| A       | (5,5)            |} \\
\texttt{| M       | (2,6)            |} \\
\\ 
\texttt{**Motivazione:**} \\
\texttt{{-} `P` e `M` sono vicini in entrambi i tabelloni: li scambio di posto per evitare sovrapposizioni.} \\
\texttt{{-} `O` e `T` sono spostati verso il centro per ridurre i movimenti (tu hai `O` a (5,4), io a (4,1)).} \\
\texttt{{-} `S` e `A` sono posizionati in celle libere e vicine tra loro (tu hai `S` a (6,2) e `A` a (6,5), io ho `S` a (2,5) e `A` a (3,5)).} \\
\texttt{{-} `D` è già a (4,4) per entrambi: lo lasciamo lì.} \\
\\ 
\texttt{\#\#\# Primo comando:} \\
\texttt{Invio un messaggio per proporre la configurazione e chiedere conferma:} \\
\texttt{```} \\
\texttt{DIRE: Proposta configurazione obiettivo: P(3,6), O(4,2), T(6,2), S(5,3), D(4,4), A(5,5), M(2,6). Confermi o hai suggerimenti?} \\
\texttt{```} \\
\\ 
\texttt{Aspetto la tua risposta per procedere con i movimenti!} \\
            }
        }
    }
     \\ \\

    \theutterance \stepcounter{utterance}  
    & & & \multicolumn{2}{p{0.3\linewidth}}{
        \cellcolor[rgb]{0.9,0.9,0.9}{
            \makecell[{{p{\linewidth}}}]{
                \texttt{\tiny{[GM$|$GM]}}
                \texttt{Grazie per la descrizione dettagliata! Prima di proporre una configurazione obiettivo comune, condivido le posizioni degli oggetti sul mio tabellone per facilitare il confronto:} \\
\\ 
\texttt{**Posizioni oggetti sul mio tabellone:**} \\
\texttt{{-} P(2,6)} \\
\texttt{{-} O(4,1)} \\
\texttt{{-} T(6,1)} \\
\texttt{{-} S(2,5) e A(3,5) (occupano celle adiacenti, ma sono oggetti separati)} \\
\texttt{{-} D(4,4)} \\
\texttt{{-} M(1,6)} \\
\\ 
\texttt{\#\#\# Analisi e proposta:} \\
\texttt{Vedo che alcuni oggetti sono già vicini tra i due tabelloni (es. `P` è a (2,6) per me e (3,5) per te, `O` è a (4,1) per me e (5,4) per te). Per minimizzare i movimenti, propongo questa **configurazione obiettivo comune**:} \\
\\ 
\texttt{| Oggetto | Coordinate (x,y) |} \\
\texttt{|{-}{-}{-}{-}{-}{-}{-}{-}{-}|{-}{-}{-}{-}{-}{-}{-}{-}{-}{-}{-}{-}{-}{-}{-}{-}{-}{-}|} \\
\texttt{| P       | (3,6)            |} \\
\texttt{| O       | (4,2)            |} \\
\texttt{| T       | (6,2)            |} \\
\texttt{| S       | (5,3)            |} \\
\texttt{| D       | (4,4)            |} \\
\texttt{| A       | (5,5)            |} \\
\texttt{| M       | (2,6)            |} \\
\\ 
\texttt{**Motivazione:**} \\
\texttt{{-} `P` e `M` sono vicini in entrambi i tabelloni: li scambio di posto per evitare sovrapposizioni.} \\
\texttt{{-} `O` e `T` sono spostati verso il centro per ridurre i movimenti (tu hai `O` a (5,4), io a (4,1)).} \\
\texttt{{-} `S` e `A` sono posizionati in celle libere e vicine tra loro (tu hai `S` a (6,2) e `A` a (6,5), io ho `S` a (2,5) e `A` a (3,5)).} \\
\texttt{{-} `D` è già a (4,4) per entrambi: lo lasciamo lì.} \\
\\ 
\texttt{\#\#\# Primo comando:} \\
\texttt{Invio un messaggio per proporre la configurazione e chiedere conferma:} \\
\texttt{```} \\
\texttt{DIRE: Proposta configurazione obiettivo: P(3,6), O(4,2), T(6,2), S(5,3), D(4,4), A(5,5), M(2,6). Confermi o hai suggerimenti?} \\
\texttt{```} \\
\\ 
\texttt{Aspetto la tua risposta per procedere con i movimenti!} \\
            }
        }
    }
    & & \\ \\

    \theutterance \stepcounter{utterance}  
    & & & \multicolumn{2}{p{0.3\linewidth}}{
        \cellcolor[rgb]{0.9,0.9,0.9}{
            \makecell[{{p{\linewidth}}}]{
                \texttt{\tiny{[GM$|$GM]}}
                \texttt{Invalid format: head and tail are not empty} \\
\texttt{head: 'Grazie per la descrizione dettagliata! Prima di proporre una configurazione obiettivo comune, condivido le posizioni degli oggetti sul mio tabellone per facilitare il confronto:} \\
\\ 
\texttt{**Posizioni oggetti sul mio tabellone:**} \\
\texttt{{-} P(2,6)} \\
\texttt{{-} O(4,1)} \\
\texttt{{-} T(6,1)} \\
\texttt{{-} S(2,5) e A(3,5) (occupano celle adiacenti, ma sono oggetti separati)} \\
\texttt{{-} D(4,4)} \\
\texttt{{-} M(1,6)} \\
\\ 
\texttt{\#\#\# Analisi e proposta:} \\
\texttt{Vedo che alcuni oggetti sono già vicini tra i due tabelloni (es. P è a (2,6) per me e (3,5) per te, O è a (4,1) per me e (5,4) per te). Per minimizzare i movimenti, propongo questa **configurazione obiettivo comune**:} \\
\\ 
\texttt{| Oggetto | Coordinate (x,y) |} \\
\texttt{|{-}{-}{-}{-}{-}{-}{-}{-}{-}|{-}{-}{-}{-}{-}{-}{-}{-}{-}{-}{-}{-}{-}{-}{-}{-}{-}{-}|} \\
\texttt{| P       | (3,6)            |} \\
\texttt{| O       | (4,2)            |} \\
\texttt{| T       | (6,2)            |} \\
\texttt{| S       | (5,3)            |} \\
\texttt{| D       | (4,4)            |} \\
\texttt{| A       | (5,5)            |} \\
\texttt{| M       | (2,6)            |} \\
\\ 
\texttt{**Motivazione:**} \\
\texttt{{-} P e M sono vicini in entrambi i tabelloni: li scambio di posto per evitare sovrapposizioni.} \\
\texttt{{-} O e T sono spostati verso il centro per ridurre i movimenti (tu hai O a (5,4), io a (4,1)).} \\
\texttt{{-} S e A sono posizionati in celle libere e vicine tra loro (tu hai S a (6,2) e A a (6,5), io ho S a (2,5) e A a (3,5)).} \\
\texttt{{-} D è già a (4,4) per entrambi: lo lasciamo lì.} \\
\\ 
\texttt{\#\#\# Primo comando:} \\
\texttt{Invio un messaggio per proporre la configurazione e chiedere conferma:} \\
\\ 
\texttt{'} \\
\texttt{tail: '} \\
\\ 
\\ 
\texttt{Aspetto la tua risposta per procedere con i movimenti!'} \\
            }
        }
    }
    & & \\ \\

    \theutterance \stepcounter{utterance}  
    & & \multicolumn{4}{p{0.6\linewidth}}{
        \cellcolor[rgb]{0.9,0.9,0.9}{
            \makecell[{{p{\linewidth}}}]{
                \texttt{\tiny{[GM$\rangle$P2]}}
                \texttt{Penalità: Il tuo messaggio non deve contenere nulla prima o dopo il comando!} \\
\\ 
\texttt{Assicurati che la tua risposta contenga solo `DIRE: <MESSAGGIO>` o `MUOVERE: <OGGETTO>, (<X>, <Y>)`, e nient'altro!} \\
\\ 
\texttt{Avete accumulato collettivamente 2 di 16 penalità.} \\
\texttt{Per favore riprova!} \\
            }
        }
    }
    & \\ \\

    \theutterance \stepcounter{utterance}  
    & & & \multicolumn{4}{p{0.6\linewidth}}{
        \cellcolor[rgb]{0.9,0.9,0.9}{
            \makecell[{{p{\linewidth}}}]{
                \texttt{\tiny{[GM$\langle$P2]}}
                \texttt{```} \\
\texttt{DIRE: Configurazione obiettivo proposta: P(3,6) O(4,2) T(6,2) S(5,3) D(4,4) A(5,5) M(2,6). Confermi?} \\
\texttt{```} \\
            }
        }
    }
     \\ \\

    \theutterance \stepcounter{utterance}  
    & & & \multicolumn{2}{p{0.3\linewidth}}{
        \cellcolor[rgb]{0.9,0.9,0.9}{
            \makecell[{{p{\linewidth}}}]{
                \texttt{\tiny{[GM$|$GM]}}
                \texttt{```} \\
\texttt{DIRE: Configurazione obiettivo proposta: P(3,6) O(4,2) T(6,2) S(5,3) D(4,4) A(5,5) M(2,6). Confermi?} \\
\texttt{```} \\
            }
        }
    }
    & & \\ \\

    \theutterance \stepcounter{utterance}  
    & & \multicolumn{4}{p{0.6\linewidth}}{
        \cellcolor[rgb]{0.9,0.9,0.9}{
            \makecell[{{p{\linewidth}}}]{
                \texttt{\tiny{[P1$\langle$GM]}}
                \texttt{Il tuo messaggio è stato trasmesso all'altro giocatore.} \\
\\ 
\\ 
\texttt{Stai giocando il round 1 di un massimo di 28.} \\
\texttt{Avete accumulato collettivamente 2 di 16 penalità.} \\
\texttt{L'altro giocatore ha inviato questo messaggio:} \\
\\ 
\\ 
\texttt{"Configurazione obiettivo proposta: P(3,6) O(4,2) T(6,2) S(5,3) D(4,4) A(5,5) M(2,6). Confermi?"} \\
\\ 
\\ 
\texttt{Qual è il tuo prossimo comando?} \\
            }
        }
    }
    & \\ \\

    \theutterance \stepcounter{utterance}  
    & \multicolumn{4}{p{0.6\linewidth}}{
        \cellcolor[rgb]{0.9,0.9,0.9}{
            \makecell[{{p{\linewidth}}}]{
                \texttt{\tiny{[P1$\rangle$GM]}}
                \texttt{```} \\
\texttt{DIRE: Confermo la configurazione obiettivo proposta. Inizio a spostare gli oggetti. Prima muovo P da (3,5) a (3,6).} \\
\texttt{```} \\
            }
        }
    }
    & & \\ \\

    \theutterance \stepcounter{utterance}  
    & & & \multicolumn{2}{p{0.3\linewidth}}{
        \cellcolor[rgb]{0.9,0.9,0.9}{
            \makecell[{{p{\linewidth}}}]{
                \texttt{\tiny{[GM$|$GM]}}
                \texttt{```} \\
\texttt{DIRE: Confermo la configurazione obiettivo proposta. Inizio a spostare gli oggetti. Prima muovo P da (3,5) a (3,6).} \\
\texttt{```} \\
            }
        }
    }
    & & \\ \\

    \theutterance \stepcounter{utterance}  
    & & \multicolumn{4}{p{0.6\linewidth}}{
        \cellcolor[rgb]{0.9,0.9,0.9}{
            \makecell[{{p{\linewidth}}}]{
                \texttt{\tiny{[GM$\rangle$P2]}}
                \texttt{Il tuo messaggio è stato trasmesso all'altro giocatore.} \\
\\ 
\\ 
\texttt{Stai giocando il round 2 di un massimo di 28.} \\
\texttt{Avete accumulato collettivamente 2 di 16 penalità.} \\
\texttt{L'altro giocatore ha inviato questo messaggio:} \\
\\ 
\\ 
\texttt{"Confermo la configurazione obiettivo proposta. Inizio a spostare gli oggetti. Prima muovo P da (3,5) a (3,6)."} \\
\\ 
\\ 
\texttt{Qual è il tuo prossimo comando?} \\
            }
        }
    }
    & \\ \\

    \theutterance \stepcounter{utterance}  
    & & & \multicolumn{4}{p{0.6\linewidth}}{
        \cellcolor[rgb]{0.9,0.9,0.9}{
            \makecell[{{p{\linewidth}}}]{
                \texttt{\tiny{[GM$\langle$P2]}}
                \texttt{Ora che la configurazione obiettivo è confermata, inizio a spostare gli oggetti sul mio tabellone per allinearli.} \\
\\ 
\texttt{**Analisi:**} \\
\texttt{{-} Il mio `P` è attualmente a **(2,6)**, ma l'obiettivo è **(3,6)**.} \\
\texttt{{-} La cella **(3,6)** è vuota (`◌`), quindi posso spostarlo lì.} \\
\\ 
\texttt{**Prossimo comando:**} \\
\texttt{```} \\
\texttt{MUOVERE: P, (3, 6)} \\
\texttt{```} \\
            }
        }
    }
     \\ \\

    \theutterance \stepcounter{utterance}  
    & & & \multicolumn{2}{p{0.3\linewidth}}{
        \cellcolor[rgb]{0.9,0.9,0.9}{
            \makecell[{{p{\linewidth}}}]{
                \texttt{\tiny{[GM$|$GM]}}
                \texttt{Ora che la configurazione obiettivo è confermata, inizio a spostare gli oggetti sul mio tabellone per allinearli.} \\
\\ 
\texttt{**Analisi:**} \\
\texttt{{-} Il mio `P` è attualmente a **(2,6)**, ma l'obiettivo è **(3,6)**.} \\
\texttt{{-} La cella **(3,6)** è vuota (`◌`), quindi posso spostarlo lì.} \\
\\ 
\texttt{**Prossimo comando:**} \\
\texttt{```} \\
\texttt{MUOVERE: P, (3, 6)} \\
\texttt{```} \\
            }
        }
    }
    & & \\ \\

    \theutterance \stepcounter{utterance}  
    & & & \multicolumn{2}{p{0.3\linewidth}}{
        \cellcolor[rgb]{0.9,0.9,0.9}{
            \makecell[{{p{\linewidth}}}]{
                \texttt{\tiny{[GM$|$GM]}}
                \texttt{Invalid format: head is not empty: 'Ora che la configurazione obiettivo è confermata, inizio a spostare gli oggetti sul mio tabellone per allinearli.} \\
\\ 
\texttt{**Analisi:**} \\
\texttt{{-} Il mio P è attualmente a **(2,6)**, ma l'obiettivo è **(3,6)**.} \\
\texttt{{-} La cella **(3,6)** è vuota (◌), quindi posso spostarlo lì.} \\
\\ 
\texttt{**Prossimo comando:**} \\
\\ 
\texttt{'} \\
            }
        }
    }
    & & \\ \\

    \theutterance \stepcounter{utterance}  
    & & \multicolumn{4}{p{0.6\linewidth}}{
        \cellcolor[rgb]{0.9,0.9,0.9}{
            \makecell[{{p{\linewidth}}}]{
                \texttt{\tiny{[GM$\rangle$P2]}}
                \texttt{Penalità: Il tuo messaggio non deve contenere nulla prima del comando!} \\
\\ 
\texttt{Assicurati che la tua risposta contenga solo `DIRE: <MESSAGGIO>` o `MUOVERE: <OGGETTO>, (<X>, <Y>)`, e nient'altro!} \\
\\ 
\texttt{Avete accumulato collettivamente 3 di 16 penalità.} \\
\texttt{Per favore riprova!} \\
            }
        }
    }
    & \\ \\

    \theutterance \stepcounter{utterance}  
    & & & \multicolumn{4}{p{0.6\linewidth}}{
        \cellcolor[rgb]{0.9,0.9,0.9}{
            \makecell[{{p{\linewidth}}}]{
                \texttt{\tiny{[GM$\langle$P2]}}
                \texttt{```} \\
\texttt{MUOVERE: P, (3, 6)} \\
\texttt{```} \\
            }
        }
    }
     \\ \\

    \theutterance \stepcounter{utterance}  
    & & & \multicolumn{2}{p{0.3\linewidth}}{
        \cellcolor[rgb]{0.9,0.9,0.9}{
            \makecell[{{p{\linewidth}}}]{
                \texttt{\tiny{[GM$|$GM]}}
                \texttt{```} \\
\texttt{MUOVERE: P, (3, 6)} \\
\texttt{```} \\
            }
        }
    }
    & & \\ \\

    \theutterance \stepcounter{utterance}  
    & & & \multicolumn{2}{p{0.3\linewidth}}{
        \cellcolor[rgb]{0.9,0.9,0.9}{
            \makecell[{{p{\linewidth}}}]{
                \texttt{\tiny{[GM$|$GM]}}
                \texttt{Penalità: (3,6) non è vuoto, ma contiene '┌'.} \\
\\ 
\texttt{Avete accumulato collettivamente 4 di 16 penalità.} \\
\texttt{Per favore riprova!} \\
            }
        }
    }
    & & \\ \\

    \theutterance \stepcounter{utterance}  
    & & \multicolumn{4}{p{0.6\linewidth}}{
        \cellcolor[rgb]{0.9,0.9,0.9}{
            \makecell[{{p{\linewidth}}}]{
                \texttt{\tiny{[GM$\rangle$P2]}}
                \texttt{Penalità: (3,6) non è vuoto, ma contiene '┌'.} \\
\\ 
\texttt{Avete accumulato collettivamente 4 di 16 penalità.} \\
\texttt{Per favore riprova!} \\
            }
        }
    }
    & \\ \\

    \theutterance \stepcounter{utterance}  
    & & & \multicolumn{4}{p{0.6\linewidth}}{
        \cellcolor[rgb]{0.9,0.9,0.9}{
            \makecell[{{p{\linewidth}}}]{
                \texttt{\tiny{[GM$\langle$P2]}}
                \texttt{Chiedo scusa per l'errore! Ho controllato male la griglia: la cella **(3,6)** contiene effettivamente un angolo (`┌`), quindi non è valida.} \\
\\ 
\texttt{**Nuovo tentativo:**} \\
\texttt{{-} Il mio `P` è a **(2,6)**.} \\
\texttt{{-} La cella **(3,5)** è vuota (`◌`), ma è già occupata da `A` nel mio tabellone (a **(3,5)**).} \\
\texttt{{-} La cella **(4,6)** è vuota (`◌`), ma è fuori dal mio obiettivo.} \\
\texttt{{-} La cella **(1,6)** è occupata da `M`.} \\
\texttt{{-} La cella **(2,5)** è vuota (`◌`), ma non è utile per l'obiettivo.} \\
\\ 
\texttt{**Soluzione:**} \\
\texttt{Sposto `P` temporaneamente in una cella vuota vicina per liberare spazio, ad esempio **(1,5)** (vuota in base alla griglia).} \\
\\ 
\texttt{```} \\
\texttt{MUOVERE: P, (1, 5)} \\
\texttt{```} \\
            }
        }
    }
     \\ \\

    \theutterance \stepcounter{utterance}  
    & & & \multicolumn{2}{p{0.3\linewidth}}{
        \cellcolor[rgb]{0.9,0.9,0.9}{
            \makecell[{{p{\linewidth}}}]{
                \texttt{\tiny{[GM$|$GM]}}
                \texttt{Chiedo scusa per l'errore! Ho controllato male la griglia: la cella **(3,6)** contiene effettivamente un angolo (`┌`), quindi non è valida.} \\
\\ 
\texttt{**Nuovo tentativo:**} \\
\texttt{{-} Il mio `P` è a **(2,6)**.} \\
\texttt{{-} La cella **(3,5)** è vuota (`◌`), ma è già occupata da `A` nel mio tabellone (a **(3,5)**).} \\
\texttt{{-} La cella **(4,6)** è vuota (`◌`), ma è fuori dal mio obiettivo.} \\
\texttt{{-} La cella **(1,6)** è occupata da `M`.} \\
\texttt{{-} La cella **(2,5)** è vuota (`◌`), ma non è utile per l'obiettivo.} \\
\\ 
\texttt{**Soluzione:**} \\
\texttt{Sposto `P` temporaneamente in una cella vuota vicina per liberare spazio, ad esempio **(1,5)** (vuota in base alla griglia).} \\
\\ 
\texttt{```} \\
\texttt{MUOVERE: P, (1, 5)} \\
\texttt{```} \\
            }
        }
    }
    & & \\ \\

    \theutterance \stepcounter{utterance}  
    & & & \multicolumn{2}{p{0.3\linewidth}}{
        \cellcolor[rgb]{0.9,0.9,0.9}{
            \makecell[{{p{\linewidth}}}]{
                \texttt{\tiny{[GM$|$GM]}}
                \texttt{Invalid format: head is not empty: 'Chiedo scusa per l'errore! Ho controllato male la griglia: la cella **(3,6)** contiene effettivamente un angolo (┌), quindi non è valida.} \\
\\ 
\texttt{**Nuovo tentativo:**} \\
\texttt{{-} Il mio P è a **(2,6)**.} \\
\texttt{{-} La cella **(3,5)** è vuota (◌), ma è già occupata da A nel mio tabellone (a **(3,5)**).} \\
\texttt{{-} La cella **(4,6)** è vuota (◌), ma è fuori dal mio obiettivo.} \\
\texttt{{-} La cella **(1,6)** è occupata da M.} \\
\texttt{{-} La cella **(2,5)** è vuota (◌), ma non è utile per l'obiettivo.} \\
\\ 
\texttt{**Soluzione:**} \\
\texttt{Sposto P temporaneamente in una cella vuota vicina per liberare spazio, ad esempio **(1,5)** (vuota in base alla griglia).} \\
\\ 
\\ 
\texttt{'} \\
            }
        }
    }
    & & \\ \\

    \theutterance \stepcounter{utterance}  
    & & \multicolumn{4}{p{0.6\linewidth}}{
        \cellcolor[rgb]{0.9,0.9,0.9}{
            \makecell[{{p{\linewidth}}}]{
                \texttt{\tiny{[GM$\rangle$P2]}}
                \texttt{Penalità: Il tuo messaggio non deve contenere nulla prima del comando!} \\
\\ 
\texttt{Assicurati che la tua risposta contenga solo `DIRE: <MESSAGGIO>` o `MUOVERE: <OGGETTO>, (<X>, <Y>)`, e nient'altro!} \\
\\ 
\texttt{Avete accumulato collettivamente 5 di 16 penalità.} \\
\texttt{Per favore riprova!} \\
            }
        }
    }
    & \\ \\

    \theutterance \stepcounter{utterance}  
    & & & \multicolumn{4}{p{0.6\linewidth}}{
        \cellcolor[rgb]{0.9,0.9,0.9}{
            \makecell[{{p{\linewidth}}}]{
                \texttt{\tiny{[GM$\langle$P2]}}
                \texttt{```} \\
\texttt{MUOVERE: P, (1, 5)} \\
\texttt{```} \\
            }
        }
    }
     \\ \\

    \theutterance \stepcounter{utterance}  
    & & & \multicolumn{2}{p{0.3\linewidth}}{
        \cellcolor[rgb]{0.9,0.9,0.9}{
            \makecell[{{p{\linewidth}}}]{
                \texttt{\tiny{[GM$|$GM]}}
                \texttt{```} \\
\texttt{MUOVERE: P, (1, 5)} \\
\texttt{```} \\
            }
        }
    }
    & & \\ \\

    \theutterance \stepcounter{utterance}  
    & & & \multicolumn{2}{p{0.3\linewidth}}{
        \cellcolor[rgb]{0.9,0.9,0.9}{
            \makecell[{{p{\linewidth}}}]{
                \texttt{\tiny{[GM$|$GM]}}
                \texttt{P moved to (1, 5)} \\
            }
        }
    }
    & & \\ \\

    \theutterance \stepcounter{utterance}  
    & & \multicolumn{4}{p{0.6\linewidth}}{
        \cellcolor[rgb]{0.9,0.9,0.9}{
            \makecell[{{p{\linewidth}}}]{
                \texttt{\tiny{[P1$\langle$GM]}}
                \texttt{Il tuo messaggio è stato trasmesso all'altro giocatore.} \\
\\ 
\\ 
\texttt{Stai giocando il round 2 di un massimo di 28.} \\
\texttt{Avete accumulato collettivamente 5 di 16 penalità.} \\
\texttt{L'altro giocatore ha spostato un oggetto sulla sua griglia.} \\
\\ 
\\ 
\texttt{Qual è il tuo prossimo comando?} \\
            }
        }
    }
    & \\ \\

    \theutterance \stepcounter{utterance}  
    & \multicolumn{4}{p{0.6\linewidth}}{
        \cellcolor[rgb]{0.9,0.9,0.9}{
            \makecell[{{p{\linewidth}}}]{
                \texttt{\tiny{[P1$\rangle$GM]}}
                \texttt{```} \\
\texttt{MUOVERE: P, (3, 6)} \\
\texttt{```} \\
            }
        }
    }
    & & \\ \\

    \theutterance \stepcounter{utterance}  
    & & & \multicolumn{2}{p{0.3\linewidth}}{
        \cellcolor[rgb]{0.9,0.9,0.9}{
            \makecell[{{p{\linewidth}}}]{
                \texttt{\tiny{[GM$|$GM]}}
                \texttt{```} \\
\texttt{MUOVERE: P, (3, 6)} \\
\texttt{```} \\
            }
        }
    }
    & & \\ \\

    \theutterance \stepcounter{utterance}  
    & & & \multicolumn{2}{p{0.3\linewidth}}{
        \cellcolor[rgb]{0.9,0.9,0.9}{
            \makecell[{{p{\linewidth}}}]{
                \texttt{\tiny{[GM$|$GM]}}
                \texttt{Penalità: (3,6) non è vuoto, ma contiene '┌'.} \\
\\ 
\texttt{Avete accumulato collettivamente 6 di 16 penalità.} \\
\texttt{Per favore riprova!} \\
            }
        }
    }
    & & \\ \\

    \theutterance \stepcounter{utterance}  
    & & \multicolumn{4}{p{0.6\linewidth}}{
        \cellcolor[rgb]{0.9,0.9,0.9}{
            \makecell[{{p{\linewidth}}}]{
                \texttt{\tiny{[P1$\langle$GM]}}
                \texttt{Penalità: (3,6) non è vuoto, ma contiene '┌'.} \\
\\ 
\texttt{Avete accumulato collettivamente 6 di 16 penalità.} \\
\texttt{Per favore riprova!} \\
            }
        }
    }
    & \\ \\

    \theutterance \stepcounter{utterance}  
    & \multicolumn{4}{p{0.6\linewidth}}{
        \cellcolor[rgb]{0.9,0.9,0.9}{
            \makecell[{{p{\linewidth}}}]{
                \texttt{\tiny{[P1$\rangle$GM]}}
                \texttt{```} \\
\texttt{MUOVERE: P, (3, 7)} \\
\texttt{```} \\
            }
        }
    }
    & & \\ \\

    \theutterance \stepcounter{utterance}  
    & & & \multicolumn{2}{p{0.3\linewidth}}{
        \cellcolor[rgb]{0.9,0.9,0.9}{
            \makecell[{{p{\linewidth}}}]{
                \texttt{\tiny{[GM$|$GM]}}
                \texttt{```} \\
\texttt{MUOVERE: P, (3, 7)} \\
\texttt{```} \\
            }
        }
    }
    & & \\ \\

    \theutterance \stepcounter{utterance}  
    & & & \multicolumn{2}{p{0.3\linewidth}}{
        \cellcolor[rgb]{0.9,0.9,0.9}{
            \makecell[{{p{\linewidth}}}]{
                \texttt{\tiny{[GM$|$GM]}}
                \texttt{Penalità: (3,7) non è vuoto, ma contiene '│'.} \\
\\ 
\texttt{Avete accumulato collettivamente 7 di 16 penalità.} \\
\texttt{Per favore riprova!} \\
            }
        }
    }
    & & \\ \\

    \theutterance \stepcounter{utterance}  
    & & \multicolumn{4}{p{0.6\linewidth}}{
        \cellcolor[rgb]{0.9,0.9,0.9}{
            \makecell[{{p{\linewidth}}}]{
                \texttt{\tiny{[P1$\langle$GM]}}
                \texttt{Penalità: (3,7) non è vuoto, ma contiene '│'.} \\
\\ 
\texttt{Avete accumulato collettivamente 7 di 16 penalità.} \\
\texttt{Per favore riprova!} \\
            }
        }
    }
    & \\ \\

    \theutterance \stepcounter{utterance}  
    & \multicolumn{4}{p{0.6\linewidth}}{
        \cellcolor[rgb]{0.9,0.9,0.9}{
            \makecell[{{p{\linewidth}}}]{
                \texttt{\tiny{[P1$\rangle$GM]}}
                \texttt{```} \\
\texttt{DIRE: Attenzione! La cella (3,6) ha un simbolo di cornice ('┌') e (3,7) ha '│'. Provo a spostare P in (2,6) se è vuota.} \\
\texttt{```} \\
            }
        }
    }
    & & \\ \\

    \theutterance \stepcounter{utterance}  
    & & & \multicolumn{2}{p{0.3\linewidth}}{
        \cellcolor[rgb]{0.9,0.9,0.9}{
            \makecell[{{p{\linewidth}}}]{
                \texttt{\tiny{[GM$|$GM]}}
                \texttt{```} \\
\texttt{DIRE: Attenzione! La cella (3,6) ha un simbolo di cornice ('┌') e (3,7) ha '│'. Provo a spostare P in (2,6) se è vuota.} \\
\texttt{```} \\
            }
        }
    }
    & & \\ \\

    \theutterance \stepcounter{utterance}  
    & & \multicolumn{4}{p{0.6\linewidth}}{
        \cellcolor[rgb]{0.9,0.9,0.9}{
            \makecell[{{p{\linewidth}}}]{
                \texttt{\tiny{[GM$\rangle$P2]}}
                \texttt{Spostato 'P' in (1,5) con successo. La tua griglia aggiornata appare così:} \\
\texttt{```} \\
\texttt{ 1234567} \\
\texttt{╔═╤══╤══╗} \\
\texttt{║◌│◌O│T◌║ 1} \\
\texttt{╟─┤◌◌│◌◌║ 2} \\
\texttt{║◌│◌◌├──╢ 3} \\
\texttt{╟─┘◌D│◌◌║ 4} \\
\texttt{║PSA◌│◌◌║ 5} \\
\texttt{║◌M┌─┼──╢ 6} \\
\texttt{║◌◌│◌│◌◌║ 7} \\
\texttt{╚══╧═╧══╝} \\
\texttt{```} \\
\texttt{Stai giocando il round 3 di un massimo di 28.} \\
\texttt{Avete accumulato collettivamente 7 di 16 penalità.} \\
\texttt{L'altro giocatore ha inviato questo messaggio:} \\
\\ 
\\ 
\texttt{"Attenzione! La cella (3,6) ha un simbolo di cornice ('┌') e (3,7) ha '│'. Provo a spostare P in (2,6) se è vuota."} \\
\\ 
\\ 
\texttt{Qual è il tuo prossimo comando?} \\
            }
        }
    }
    & \\ \\

    \theutterance \stepcounter{utterance}  
    & & & \multicolumn{4}{p{0.6\linewidth}}{
        \cellcolor[rgb]{0.9,0.9,0.9}{
            \makecell[{{p{\linewidth}}}]{
                \texttt{\tiny{[GM$\langle$P2]}}
                \texttt{```} \\
\texttt{DIRE: (2,6) è occupato da 'M'. Provo a spostare M in (1,6) per liberare spazio. Confermi?} \\
\texttt{```} \\
            }
        }
    }
     \\ \\

    \theutterance \stepcounter{utterance}  
    & & & \multicolumn{2}{p{0.3\linewidth}}{
        \cellcolor[rgb]{0.9,0.9,0.9}{
            \makecell[{{p{\linewidth}}}]{
                \texttt{\tiny{[GM$|$GM]}}
                \texttt{```} \\
\texttt{DIRE: (2,6) è occupato da 'M'. Provo a spostare M in (1,6) per liberare spazio. Confermi?} \\
\texttt{```} \\
            }
        }
    }
    & & \\ \\

    \theutterance \stepcounter{utterance}  
    & & \multicolumn{4}{p{0.6\linewidth}}{
        \cellcolor[rgb]{0.9,0.9,0.9}{
            \makecell[{{p{\linewidth}}}]{
                \texttt{\tiny{[P1$\langle$GM]}}
                \texttt{Il tuo messaggio è stato trasmesso all'altro giocatore.} \\
\\ 
\\ 
\texttt{Stai giocando il round 3 di un massimo di 28.} \\
\texttt{Avete accumulato collettivamente 7 di 16 penalità.} \\
\texttt{L'altro giocatore ha inviato questo messaggio:} \\
\\ 
\\ 
\texttt{"(2,6) è occupato da 'M'. Provo a spostare M in (1,6) per liberare spazio. Confermi?"} \\
\\ 
\\ 
\texttt{Qual è il tuo prossimo comando?} \\
            }
        }
    }
    & \\ \\

    \theutterance \stepcounter{utterance}  
    & \multicolumn{4}{p{0.6\linewidth}}{
        \cellcolor[rgb]{0.9,0.9,0.9}{
            \makecell[{{p{\linewidth}}}]{
                \texttt{\tiny{[P1$\rangle$GM]}}
                \texttt{```} \\
\texttt{MUOVERE: M, (1, 6)} \\
\texttt{```} \\
            }
        }
    }
    & & \\ \\

    \theutterance \stepcounter{utterance}  
    & & & \multicolumn{2}{p{0.3\linewidth}}{
        \cellcolor[rgb]{0.9,0.9,0.9}{
            \makecell[{{p{\linewidth}}}]{
                \texttt{\tiny{[GM$|$GM]}}
                \texttt{```} \\
\texttt{MUOVERE: M, (1, 6)} \\
\texttt{```} \\
            }
        }
    }
    & & \\ \\

    \theutterance \stepcounter{utterance}  
    & & & \multicolumn{2}{p{0.3\linewidth}}{
        \cellcolor[rgb]{0.9,0.9,0.9}{
            \makecell[{{p{\linewidth}}}]{
                \texttt{\tiny{[GM$|$GM]}}
                \texttt{M moved to (1, 6)} \\
            }
        }
    }
    & & \\ \\

    \theutterance \stepcounter{utterance}  
    & & \multicolumn{4}{p{0.6\linewidth}}{
        \cellcolor[rgb]{0.9,0.9,0.9}{
            \makecell[{{p{\linewidth}}}]{
                \texttt{\tiny{[GM$\rangle$P2]}}
                \texttt{Il tuo messaggio è stato trasmesso all'altro giocatore.} \\
\\ 
\\ 
\texttt{Stai giocando il round 4 di un massimo di 28.} \\
\texttt{Avete accumulato collettivamente 7 di 16 penalità.} \\
\texttt{L'altro giocatore ha spostato un oggetto sulla sua griglia.} \\
\\ 
\\ 
\texttt{Qual è il tuo prossimo comando?} \\
            }
        }
    }
    & \\ \\

    \theutterance \stepcounter{utterance}  
    & & & \multicolumn{4}{p{0.6\linewidth}}{
        \cellcolor[rgb]{0.9,0.9,0.9}{
            \makecell[{{p{\linewidth}}}]{
                \texttt{\tiny{[GM$\langle$P2]}}
                \texttt{Ora sposto **M** da **(1,6)** a **(2,6)** per liberare **(1,6)** e avvicinarmi all'obiettivo finale (M deve andare in **(2,6)**).} \\
\\ 
\texttt{```} \\
\texttt{MUOVERE: M, (2, 6)} \\
\texttt{```} \\
            }
        }
    }
     \\ \\

    \theutterance \stepcounter{utterance}  
    & & & \multicolumn{2}{p{0.3\linewidth}}{
        \cellcolor[rgb]{0.9,0.9,0.9}{
            \makecell[{{p{\linewidth}}}]{
                \texttt{\tiny{[GM$|$GM]}}
                \texttt{Ora sposto **M** da **(1,6)** a **(2,6)** per liberare **(1,6)** e avvicinarmi all'obiettivo finale (M deve andare in **(2,6)**).} \\
\\ 
\texttt{```} \\
\texttt{MUOVERE: M, (2, 6)} \\
\texttt{```} \\
            }
        }
    }
    & & \\ \\

    \theutterance \stepcounter{utterance}  
    & & & \multicolumn{2}{p{0.3\linewidth}}{
        \cellcolor[rgb]{0.9,0.9,0.9}{
            \makecell[{{p{\linewidth}}}]{
                \texttt{\tiny{[GM$|$GM]}}
                \texttt{Invalid format: head is not empty: 'Ora sposto **M** da **(1,6)** a **(2,6)** per liberare **(1,6)** e avvicinarmi all'obiettivo finale (M deve andare in **(2,6)**).} \\
\\ 
\\ 
\texttt{'} \\
            }
        }
    }
    & & \\ \\

    \theutterance \stepcounter{utterance}  
    & & \multicolumn{4}{p{0.6\linewidth}}{
        \cellcolor[rgb]{0.9,0.9,0.9}{
            \makecell[{{p{\linewidth}}}]{
                \texttt{\tiny{[GM$\rangle$P2]}}
                \texttt{Penalità: Il tuo messaggio non deve contenere nulla prima del comando!} \\
\\ 
\texttt{Assicurati che la tua risposta contenga solo `DIRE: <MESSAGGIO>` o `MUOVERE: <OGGETTO>, (<X>, <Y>)`, e nient'altro!} \\
\\ 
\texttt{Avete accumulato collettivamente 8 di 16 penalità.} \\
\texttt{Per favore riprova!} \\
            }
        }
    }
    & \\ \\

    \theutterance \stepcounter{utterance}  
    & & & \multicolumn{4}{p{0.6\linewidth}}{
        \cellcolor[rgb]{0.9,0.9,0.9}{
            \makecell[{{p{\linewidth}}}]{
                \texttt{\tiny{[GM$\langle$P2]}}
                \texttt{```} \\
\texttt{MUOVERE: M, (2, 6)} \\
\texttt{```} \\
            }
        }
    }
     \\ \\

    \theutterance \stepcounter{utterance}  
    & & & \multicolumn{2}{p{0.3\linewidth}}{
        \cellcolor[rgb]{0.9,0.9,0.9}{
            \makecell[{{p{\linewidth}}}]{
                \texttt{\tiny{[GM$|$GM]}}
                \texttt{```} \\
\texttt{MUOVERE: M, (2, 6)} \\
\texttt{```} \\
            }
        }
    }
    & & \\ \\

    \theutterance \stepcounter{utterance}  
    & & & \multicolumn{2}{p{0.3\linewidth}}{
        \cellcolor[rgb]{0.9,0.9,0.9}{
            \makecell[{{p{\linewidth}}}]{
                \texttt{\tiny{[GM$|$GM]}}
                \texttt{Penalità: (2,6) non è vuoto, ma contiene 'M'.} \\
\\ 
\texttt{Avete accumulato collettivamente 9 di 16 penalità.} \\
\texttt{Per favore riprova!} \\
            }
        }
    }
    & & \\ \\

    \theutterance \stepcounter{utterance}  
    & & \multicolumn{4}{p{0.6\linewidth}}{
        \cellcolor[rgb]{0.9,0.9,0.9}{
            \makecell[{{p{\linewidth}}}]{
                \texttt{\tiny{[GM$\rangle$P2]}}
                \texttt{Penalità: (2,6) non è vuoto, ma contiene 'M'.} \\
\\ 
\texttt{Avete accumulato collettivamente 9 di 16 penalità.} \\
\texttt{Per favore riprova!} \\
            }
        }
    }
    & \\ \\

    \theutterance \stepcounter{utterance}  
    & & & \multicolumn{4}{p{0.6\linewidth}}{
        \cellcolor[rgb]{0.9,0.9,0.9}{
            \makecell[{{p{\linewidth}}}]{
                \texttt{\tiny{[GM$\langle$P2]}}
                \texttt{Mi scuso per l'errore. La cella **(2,6)** è già occupata da `M`, quindi non posso spostarlo lì.} \\
\\ 
\texttt{**Nuovo piano:**} \\
\texttt{{-} Sposto `M` temporaneamente in **(1,7)** (vuota, come da griglia).} \\
\texttt{{-} Poi libererò spazio per `P` in **(3,6)**.} \\
\\ 
\texttt{```} \\
\texttt{MUOVERE: M, (1, 7)} \\
\texttt{```} \\
            }
        }
    }
     \\ \\

    \theutterance \stepcounter{utterance}  
    & & & \multicolumn{2}{p{0.3\linewidth}}{
        \cellcolor[rgb]{0.9,0.9,0.9}{
            \makecell[{{p{\linewidth}}}]{
                \texttt{\tiny{[GM$|$GM]}}
                \texttt{Mi scuso per l'errore. La cella **(2,6)** è già occupata da `M`, quindi non posso spostarlo lì.} \\
\\ 
\texttt{**Nuovo piano:**} \\
\texttt{{-} Sposto `M` temporaneamente in **(1,7)** (vuota, come da griglia).} \\
\texttt{{-} Poi libererò spazio per `P` in **(3,6)**.} \\
\\ 
\texttt{```} \\
\texttt{MUOVERE: M, (1, 7)} \\
\texttt{```} \\
            }
        }
    }
    & & \\ \\

    \theutterance \stepcounter{utterance}  
    & & & \multicolumn{2}{p{0.3\linewidth}}{
        \cellcolor[rgb]{0.9,0.9,0.9}{
            \makecell[{{p{\linewidth}}}]{
                \texttt{\tiny{[GM$|$GM]}}
                \texttt{Invalid format: head is not empty: 'Mi scuso per l'errore. La cella **(2,6)** è già occupata da M, quindi non posso spostarlo lì.} \\
\\ 
\texttt{**Nuovo piano:**} \\
\texttt{{-} Sposto M temporaneamente in **(1,7)** (vuota, come da griglia).} \\
\texttt{{-} Poi libererò spazio per P in **(3,6)**.} \\
\\ 
\\ 
\texttt{'} \\
            }
        }
    }
    & & \\ \\

    \theutterance \stepcounter{utterance}  
    & & \multicolumn{4}{p{0.6\linewidth}}{
        \cellcolor[rgb]{0.9,0.9,0.9}{
            \makecell[{{p{\linewidth}}}]{
                \texttt{\tiny{[GM$\rangle$P2]}}
                \texttt{Penalità: Il tuo messaggio non deve contenere nulla prima del comando!} \\
\\ 
\texttt{Assicurati che la tua risposta contenga solo `DIRE: <MESSAGGIO>` o `MUOVERE: <OGGETTO>, (<X>, <Y>)`, e nient'altro!} \\
\\ 
\texttt{Avete accumulato collettivamente 10 di 16 penalità.} \\
\texttt{Per favore riprova!} \\
            }
        }
    }
    & \\ \\

    \theutterance \stepcounter{utterance}  
    & & & \multicolumn{4}{p{0.6\linewidth}}{
        \cellcolor[rgb]{0.9,0.9,0.9}{
            \makecell[{{p{\linewidth}}}]{
                \texttt{\tiny{[GM$\langle$P2]}}
                \texttt{```} \\
\texttt{MUOVERE: M, (1, 7)} \\
\texttt{```} \\
            }
        }
    }
     \\ \\

    \theutterance \stepcounter{utterance}  
    & & & \multicolumn{2}{p{0.3\linewidth}}{
        \cellcolor[rgb]{0.9,0.9,0.9}{
            \makecell[{{p{\linewidth}}}]{
                \texttt{\tiny{[GM$|$GM]}}
                \texttt{```} \\
\texttt{MUOVERE: M, (1, 7)} \\
\texttt{```} \\
            }
        }
    }
    & & \\ \\

    \theutterance \stepcounter{utterance}  
    & & & \multicolumn{2}{p{0.3\linewidth}}{
        \cellcolor[rgb]{0.9,0.9,0.9}{
            \makecell[{{p{\linewidth}}}]{
                \texttt{\tiny{[GM$|$GM]}}
                \texttt{M moved to (1, 7)} \\
            }
        }
    }
    & & \\ \\

    \theutterance \stepcounter{utterance}  
    & & \multicolumn{4}{p{0.6\linewidth}}{
        \cellcolor[rgb]{0.9,0.9,0.9}{
            \makecell[{{p{\linewidth}}}]{
                \texttt{\tiny{[P1$\langle$GM]}}
                \texttt{Spostato 'M' in (1,6) con successo. La tua griglia aggiornata appare così:} \\
\texttt{```} \\
\texttt{ 1234567} \\
\texttt{╔═╤══╤══╗} \\
\texttt{║◌│◌◌│◌◌║ 1} \\
\texttt{╟─┤◌◌│◌S║ 2} \\
\texttt{║◌│TD├──╢ 3} \\
\texttt{╟─┘◌◌│O◌║ 4} \\
\texttt{║◌◌P◌│◌A║ 5} \\
\texttt{║M◌┌─┼──╢ 6} \\
\texttt{║◌◌│◌│◌◌║ 7} \\
\texttt{╚══╧═╧══╝} \\
\texttt{```} \\
\texttt{Stai giocando il round 4 di un massimo di 28.} \\
\texttt{Avete accumulato collettivamente 10 di 16 penalità.} \\
\texttt{L'altro giocatore ha spostato un oggetto sulla sua griglia.} \\
\\ 
\\ 
\texttt{Qual è il tuo prossimo comando?} \\
            }
        }
    }
    & \\ \\

    \theutterance \stepcounter{utterance}  
    & \multicolumn{4}{p{0.6\linewidth}}{
        \cellcolor[rgb]{0.9,0.9,0.9}{
            \makecell[{{p{\linewidth}}}]{
                \texttt{\tiny{[P1$\rangle$GM]}}
                \texttt{```} \\
\texttt{MUOVERE: P, (3, 6)} \\
\texttt{```} \\
            }
        }
    }
    & & \\ \\

    \theutterance \stepcounter{utterance}  
    & & & \multicolumn{2}{p{0.3\linewidth}}{
        \cellcolor[rgb]{0.9,0.9,0.9}{
            \makecell[{{p{\linewidth}}}]{
                \texttt{\tiny{[GM$|$GM]}}
                \texttt{```} \\
\texttt{MUOVERE: P, (3, 6)} \\
\texttt{```} \\
            }
        }
    }
    & & \\ \\

    \theutterance \stepcounter{utterance}  
    & & & \multicolumn{2}{p{0.3\linewidth}}{
        \cellcolor[rgb]{0.9,0.9,0.9}{
            \makecell[{{p{\linewidth}}}]{
                \texttt{\tiny{[GM$|$GM]}}
                \texttt{Penalità: (3,6) non è vuoto, ma contiene '┌'.} \\
\\ 
\texttt{Avete accumulato collettivamente 11 di 16 penalità.} \\
\texttt{Per favore riprova!} \\
            }
        }
    }
    & & \\ \\

    \theutterance \stepcounter{utterance}  
    & & \multicolumn{4}{p{0.6\linewidth}}{
        \cellcolor[rgb]{0.9,0.9,0.9}{
            \makecell[{{p{\linewidth}}}]{
                \texttt{\tiny{[P1$\langle$GM]}}
                \texttt{Penalità: (3,6) non è vuoto, ma contiene '┌'.} \\
\\ 
\texttt{Avete accumulato collettivamente 11 di 16 penalità.} \\
\texttt{Per favore riprova!} \\
            }
        }
    }
    & \\ \\

    \theutterance \stepcounter{utterance}  
    & \multicolumn{4}{p{0.6\linewidth}}{
        \cellcolor[rgb]{0.9,0.9,0.9}{
            \makecell[{{p{\linewidth}}}]{
                \texttt{\tiny{[P1$\rangle$GM]}}
                \texttt{```} \\
\texttt{DIRE: La cella (3,6) è ancora occupata da '┌'. Provo a spostare P in (4,6) se è libera. Confermi che è vuota?} \\
\texttt{```} \\
            }
        }
    }
    & & \\ \\

    \theutterance \stepcounter{utterance}  
    & & & \multicolumn{2}{p{0.3\linewidth}}{
        \cellcolor[rgb]{0.9,0.9,0.9}{
            \makecell[{{p{\linewidth}}}]{
                \texttt{\tiny{[GM$|$GM]}}
                \texttt{```} \\
\texttt{DIRE: La cella (3,6) è ancora occupata da '┌'. Provo a spostare P in (4,6) se è libera. Confermi che è vuota?} \\
\texttt{```} \\
            }
        }
    }
    & & \\ \\

    \theutterance \stepcounter{utterance}  
    & & \multicolumn{4}{p{0.6\linewidth}}{
        \cellcolor[rgb]{0.9,0.9,0.9}{
            \makecell[{{p{\linewidth}}}]{
                \texttt{\tiny{[GM$\rangle$P2]}}
                \texttt{Spostato 'M' in (1,7) con successo. La tua griglia aggiornata appare così:} \\
\texttt{```} \\
\texttt{ 1234567} \\
\texttt{╔═╤══╤══╗} \\
\texttt{║◌│◌O│T◌║ 1} \\
\texttt{╟─┤◌◌│◌◌║ 2} \\
\texttt{║◌│◌◌├──╢ 3} \\
\texttt{╟─┘◌D│◌◌║ 4} \\
\texttt{║PSA◌│◌◌║ 5} \\
\texttt{║◌◌┌─┼──╢ 6} \\
\texttt{║M◌│◌│◌◌║ 7} \\
\texttt{╚══╧═╧══╝} \\
\texttt{```} \\
\texttt{Stai giocando il round 5 di un massimo di 28.} \\
\texttt{Avete accumulato collettivamente 11 di 16 penalità.} \\
\texttt{L'altro giocatore ha inviato questo messaggio:} \\
\\ 
\\ 
\texttt{"La cella (3,6) è ancora occupata da '┌'. Provo a spostare P in (4,6) se è libera. Confermi che è vuota?"} \\
\\ 
\\ 
\texttt{Qual è il tuo prossimo comando?} \\
            }
        }
    }
    & \\ \\

    \theutterance \stepcounter{utterance}  
    & & & \multicolumn{4}{p{0.6\linewidth}}{
        \cellcolor[rgb]{0.9,0.9,0.9}{
            \makecell[{{p{\linewidth}}}]{
                \texttt{\tiny{[GM$\langle$P2]}}
                \texttt{```} \\
\texttt{DIRE: (4,6) è occupata da '─'. La cella libera più vicina a (3,6) è (4,5) (vuota). Vuoi che sposti P lì?} \\
\texttt{```} \\
            }
        }
    }
     \\ \\

    \theutterance \stepcounter{utterance}  
    & & & \multicolumn{2}{p{0.3\linewidth}}{
        \cellcolor[rgb]{0.9,0.9,0.9}{
            \makecell[{{p{\linewidth}}}]{
                \texttt{\tiny{[GM$|$GM]}}
                \texttt{```} \\
\texttt{DIRE: (4,6) è occupata da '─'. La cella libera più vicina a (3,6) è (4,5) (vuota). Vuoi che sposti P lì?} \\
\texttt{```} \\
            }
        }
    }
    & & \\ \\

    \theutterance \stepcounter{utterance}  
    & & \multicolumn{4}{p{0.6\linewidth}}{
        \cellcolor[rgb]{0.9,0.9,0.9}{
            \makecell[{{p{\linewidth}}}]{
                \texttt{\tiny{[P1$\langle$GM]}}
                \texttt{Il tuo messaggio è stato trasmesso all'altro giocatore.} \\
\\ 
\\ 
\texttt{Stai giocando il round 5 di un massimo di 28.} \\
\texttt{Avete accumulato collettivamente 11 di 16 penalità.} \\
\texttt{L'altro giocatore ha inviato questo messaggio:} \\
\\ 
\\ 
\texttt{"(4,6) è occupata da '─'. La cella libera più vicina a (3,6) è (4,5) (vuota). Vuoi che sposti P lì?"} \\
\\ 
\\ 
\texttt{Qual è il tuo prossimo comando?} \\
            }
        }
    }
    & \\ \\

    \theutterance \stepcounter{utterance}  
    & \multicolumn{4}{p{0.6\linewidth}}{
        \cellcolor[rgb]{0.9,0.9,0.9}{
            \makecell[{{p{\linewidth}}}]{
                \texttt{\tiny{[P1$\rangle$GM]}}
                \texttt{```} \\
\texttt{MUOVERE: P, (4, 5)} \\
\texttt{```} \\
            }
        }
    }
    & & \\ \\

    \theutterance \stepcounter{utterance}  
    & & & \multicolumn{2}{p{0.3\linewidth}}{
        \cellcolor[rgb]{0.9,0.9,0.9}{
            \makecell[{{p{\linewidth}}}]{
                \texttt{\tiny{[GM$|$GM]}}
                \texttt{```} \\
\texttt{MUOVERE: P, (4, 5)} \\
\texttt{```} \\
            }
        }
    }
    & & \\ \\

    \theutterance \stepcounter{utterance}  
    & & & \multicolumn{2}{p{0.3\linewidth}}{
        \cellcolor[rgb]{0.9,0.9,0.9}{
            \makecell[{{p{\linewidth}}}]{
                \texttt{\tiny{[GM$|$GM]}}
                \texttt{P moved to (4, 5)} \\
            }
        }
    }
    & & \\ \\

    \theutterance \stepcounter{utterance}  
    & & \multicolumn{4}{p{0.6\linewidth}}{
        \cellcolor[rgb]{0.9,0.9,0.9}{
            \makecell[{{p{\linewidth}}}]{
                \texttt{\tiny{[GM$\rangle$P2]}}
                \texttt{Il tuo messaggio è stato trasmesso all'altro giocatore.} \\
\\ 
\\ 
\texttt{Stai giocando il round 6 di un massimo di 28.} \\
\texttt{Avete accumulato collettivamente 11 di 16 penalità.} \\
\texttt{L'altro giocatore ha spostato un oggetto sulla sua griglia.} \\
\\ 
\\ 
\texttt{Qual è il tuo prossimo comando?} \\
            }
        }
    }
    & \\ \\

    \theutterance \stepcounter{utterance}  
    & & & \multicolumn{4}{p{0.6\linewidth}}{
        \cellcolor[rgb]{0.9,0.9,0.9}{
            \makecell[{{p{\linewidth}}}]{
                \texttt{\tiny{[GM$\langle$P2]}}
                \texttt{Ora che ho liberato (1,6) spostando M in (1,7), posso finalmente muovere P verso la sua destinazione finale (3,6). Prima però devo liberare la strada:} \\
\\ 
\texttt{1. Sposto **S** da (2,5) a (1,5) per liberare (2,5)} \\
\texttt{2. Poi sposto **A** da (3,5) a (2,5)} \\
\texttt{3. Infine sposto **P** da (1,5) a (3,6)} \\
\\ 
\texttt{**Primo passo:**} \\
\texttt{```} \\
\texttt{MUOVERE: S, (1, 5)} \\
\texttt{```} \\
            }
        }
    }
     \\ \\

    \theutterance \stepcounter{utterance}  
    & & & \multicolumn{2}{p{0.3\linewidth}}{
        \cellcolor[rgb]{0.9,0.9,0.9}{
            \makecell[{{p{\linewidth}}}]{
                \texttt{\tiny{[GM$|$GM]}}
                \texttt{Ora che ho liberato (1,6) spostando M in (1,7), posso finalmente muovere P verso la sua destinazione finale (3,6). Prima però devo liberare la strada:} \\
\\ 
\texttt{1. Sposto **S** da (2,5) a (1,5) per liberare (2,5)} \\
\texttt{2. Poi sposto **A** da (3,5) a (2,5)} \\
\texttt{3. Infine sposto **P** da (1,5) a (3,6)} \\
\\ 
\texttt{**Primo passo:**} \\
\texttt{```} \\
\texttt{MUOVERE: S, (1, 5)} \\
\texttt{```} \\
            }
        }
    }
    & & \\ \\

    \theutterance \stepcounter{utterance}  
    & & & \multicolumn{2}{p{0.3\linewidth}}{
        \cellcolor[rgb]{0.9,0.9,0.9}{
            \makecell[{{p{\linewidth}}}]{
                \texttt{\tiny{[GM$|$GM]}}
                \texttt{Invalid format: head is not empty: 'Ora che ho liberato (1,6) spostando M in (1,7), posso finalmente muovere P verso la sua destinazione finale (3,6). Prima però devo liberare la strada:} \\
\\ 
\texttt{1. Sposto **S** da (2,5) a (1,5) per liberare (2,5)} \\
\texttt{2. Poi sposto **A** da (3,5) a (2,5)} \\
\texttt{3. Infine sposto **P** da (1,5) a (3,6)} \\
\\ 
\texttt{**Primo passo:**} \\
\\ 
\texttt{'} \\
            }
        }
    }
    & & \\ \\

    \theutterance \stepcounter{utterance}  
    & & \multicolumn{4}{p{0.6\linewidth}}{
        \cellcolor[rgb]{0.9,0.9,0.9}{
            \makecell[{{p{\linewidth}}}]{
                \texttt{\tiny{[GM$\rangle$P2]}}
                \texttt{Penalità: Il tuo messaggio non deve contenere nulla prima del comando!} \\
\\ 
\texttt{Assicurati che la tua risposta contenga solo `DIRE: <MESSAGGIO>` o `MUOVERE: <OGGETTO>, (<X>, <Y>)`, e nient'altro!} \\
\\ 
\texttt{Avete accumulato collettivamente 12 di 16 penalità.} \\
\texttt{Per favore riprova!} \\
            }
        }
    }
    & \\ \\

    \theutterance \stepcounter{utterance}  
    & & & \multicolumn{4}{p{0.6\linewidth}}{
        \cellcolor[rgb]{0.9,0.9,0.9}{
            \makecell[{{p{\linewidth}}}]{
                \texttt{\tiny{[GM$\langle$P2]}}
                \texttt{```} \\
\texttt{MUOVERE: S, (1, 5)} \\
\texttt{```} \\
            }
        }
    }
     \\ \\

    \theutterance \stepcounter{utterance}  
    & & & \multicolumn{2}{p{0.3\linewidth}}{
        \cellcolor[rgb]{0.9,0.9,0.9}{
            \makecell[{{p{\linewidth}}}]{
                \texttt{\tiny{[GM$|$GM]}}
                \texttt{```} \\
\texttt{MUOVERE: S, (1, 5)} \\
\texttt{```} \\
            }
        }
    }
    & & \\ \\

    \theutterance \stepcounter{utterance}  
    & & & \multicolumn{2}{p{0.3\linewidth}}{
        \cellcolor[rgb]{0.9,0.9,0.9}{
            \makecell[{{p{\linewidth}}}]{
                \texttt{\tiny{[GM$|$GM]}}
                \texttt{Penalità: (1,5) non è vuoto, ma contiene 'P'.} \\
\\ 
\texttt{Avete accumulato collettivamente 13 di 16 penalità.} \\
\texttt{Per favore riprova!} \\
            }
        }
    }
    & & \\ \\

    \theutterance \stepcounter{utterance}  
    & & \multicolumn{4}{p{0.6\linewidth}}{
        \cellcolor[rgb]{0.9,0.9,0.9}{
            \makecell[{{p{\linewidth}}}]{
                \texttt{\tiny{[GM$\rangle$P2]}}
                \texttt{Penalità: (1,5) non è vuoto, ma contiene 'P'.} \\
\\ 
\texttt{Avete accumulato collettivamente 13 di 16 penalità.} \\
\texttt{Per favore riprova!} \\
            }
        }
    }
    & \\ \\

    \theutterance \stepcounter{utterance}  
    & & & \multicolumn{4}{p{0.6\linewidth}}{
        \cellcolor[rgb]{0.9,0.9,0.9}{
            \makecell[{{p{\linewidth}}}]{
                \texttt{\tiny{[GM$\langle$P2]}}
                \texttt{```} \\
\texttt{MUOVERE: S, (7, 5)} \\
\texttt{```} \\
            }
        }
    }
     \\ \\

    \theutterance \stepcounter{utterance}  
    & & & \multicolumn{2}{p{0.3\linewidth}}{
        \cellcolor[rgb]{0.9,0.9,0.9}{
            \makecell[{{p{\linewidth}}}]{
                \texttt{\tiny{[GM$|$GM]}}
                \texttt{```} \\
\texttt{MUOVERE: S, (7, 5)} \\
\texttt{```} \\
            }
        }
    }
    & & \\ \\

    \theutterance \stepcounter{utterance}  
    & & & \multicolumn{2}{p{0.3\linewidth}}{
        \cellcolor[rgb]{0.9,0.9,0.9}{
            \makecell[{{p{\linewidth}}}]{
                \texttt{\tiny{[GM$|$GM]}}
                \texttt{S moved to (7, 5)} \\
            }
        }
    }
    & & \\ \\

    \theutterance \stepcounter{utterance}  
    & & \multicolumn{4}{p{0.6\linewidth}}{
        \cellcolor[rgb]{0.9,0.9,0.9}{
            \makecell[{{p{\linewidth}}}]{
                \texttt{\tiny{[P1$\langle$GM]}}
                \texttt{Spostato 'P' in (4,5) con successo. La tua griglia aggiornata appare così:} \\
\texttt{```} \\
\texttt{ 1234567} \\
\texttt{╔═╤══╤══╗} \\
\texttt{║◌│◌◌│◌◌║ 1} \\
\texttt{╟─┤◌◌│◌S║ 2} \\
\texttt{║◌│TD├──╢ 3} \\
\texttt{╟─┘◌◌│O◌║ 4} \\
\texttt{║◌◌◌P│◌A║ 5} \\
\texttt{║M◌┌─┼──╢ 6} \\
\texttt{║◌◌│◌│◌◌║ 7} \\
\texttt{╚══╧═╧══╝} \\
\texttt{```} \\
\texttt{Stai giocando il round 6 di un massimo di 28.} \\
\texttt{Avete accumulato collettivamente 13 di 16 penalità.} \\
\texttt{L'altro giocatore ha spostato un oggetto sulla sua griglia.} \\
\\ 
\\ 
\texttt{Qual è il tuo prossimo comando?} \\
            }
        }
    }
    & \\ \\

    \theutterance \stepcounter{utterance}  
    & \multicolumn{4}{p{0.6\linewidth}}{
        \cellcolor[rgb]{0.9,0.9,0.9}{
            \makecell[{{p{\linewidth}}}]{
                \texttt{\tiny{[P1$\rangle$GM]}}
                \texttt{```} \\
\texttt{DIRE: P è ora in (4,5). Prossimo obiettivo: spostare O da (5,4) a (4,2). Prima devo liberare (4,2). La cella è occupata da '◌' (vuota) o c'è altro?} \\
\texttt{```} \\
            }
        }
    }
    & & \\ \\

    \theutterance \stepcounter{utterance}  
    & & & \multicolumn{2}{p{0.3\linewidth}}{
        \cellcolor[rgb]{0.9,0.9,0.9}{
            \makecell[{{p{\linewidth}}}]{
                \texttt{\tiny{[GM$|$GM]}}
                \texttt{```} \\
\texttt{DIRE: P è ora in (4,5). Prossimo obiettivo: spostare O da (5,4) a (4,2). Prima devo liberare (4,2). La cella è occupata da '◌' (vuota) o c'è altro?} \\
\texttt{```} \\
            }
        }
    }
    & & \\ \\

    \theutterance \stepcounter{utterance}  
    & & \multicolumn{4}{p{0.6\linewidth}}{
        \cellcolor[rgb]{0.9,0.9,0.9}{
            \makecell[{{p{\linewidth}}}]{
                \texttt{\tiny{[GM$\rangle$P2]}}
                \texttt{Spostato 'S' in (7,5) con successo. La tua griglia aggiornata appare così:} \\
\texttt{```} \\
\texttt{ 1234567} \\
\texttt{╔═╤══╤══╗} \\
\texttt{║◌│◌O│T◌║ 1} \\
\texttt{╟─┤◌◌│◌◌║ 2} \\
\texttt{║◌│◌◌├──╢ 3} \\
\texttt{╟─┘◌D│◌◌║ 4} \\
\texttt{║P◌A◌│◌S║ 5} \\
\texttt{║◌◌┌─┼──╢ 6} \\
\texttt{║M◌│◌│◌◌║ 7} \\
\texttt{╚══╧═╧══╝} \\
\texttt{```} \\
\texttt{Stai giocando il round 7 di un massimo di 28.} \\
\texttt{Avete accumulato collettivamente 13 di 16 penalità.} \\
\texttt{L'altro giocatore ha inviato questo messaggio:} \\
\\ 
\\ 
\texttt{"P è ora in (4,5). Prossimo obiettivo: spostare O da (5,4) a (4,2). Prima devo liberare (4,2). La cella è occupata da '◌' (vuota) o c'è altro?"} \\
\\ 
\\ 
\texttt{Qual è il tuo prossimo comando?} \\
            }
        }
    }
    & \\ \\

    \theutterance \stepcounter{utterance}  
    & & & \multicolumn{4}{p{0.6\linewidth}}{
        \cellcolor[rgb]{0.9,0.9,0.9}{
            \makecell[{{p{\linewidth}}}]{
                \texttt{\tiny{[GM$\langle$P2]}}
                \texttt{```} \\
\texttt{DIRE: (4,2) sul mio tabellone è vuota (◌). Posso spostare O da (4,1) a (4,2) subito.} \\
\texttt{```} \\
            }
        }
    }
     \\ \\

    \theutterance \stepcounter{utterance}  
    & & & \multicolumn{2}{p{0.3\linewidth}}{
        \cellcolor[rgb]{0.9,0.9,0.9}{
            \makecell[{{p{\linewidth}}}]{
                \texttt{\tiny{[GM$|$GM]}}
                \texttt{```} \\
\texttt{DIRE: (4,2) sul mio tabellone è vuota (◌). Posso spostare O da (4,1) a (4,2) subito.} \\
\texttt{```} \\
            }
        }
    }
    & & \\ \\

    \theutterance \stepcounter{utterance}  
    & & \multicolumn{4}{p{0.6\linewidth}}{
        \cellcolor[rgb]{0.9,0.9,0.9}{
            \makecell[{{p{\linewidth}}}]{
                \texttt{\tiny{[P1$\langle$GM]}}
                \texttt{Il tuo messaggio è stato trasmesso all'altro giocatore.} \\
\\ 
\\ 
\texttt{Stai giocando il round 7 di un massimo di 28.} \\
\texttt{Avete accumulato collettivamente 13 di 16 penalità.} \\
\texttt{L'altro giocatore ha inviato questo messaggio:} \\
\\ 
\\ 
\texttt{"(4,2) sul mio tabellone è vuota (◌). Posso spostare O da (4,1) a (4,2) subito."} \\
\\ 
\\ 
\texttt{Qual è il tuo prossimo comando?} \\
            }
        }
    }
    & \\ \\

    \theutterance \stepcounter{utterance}  
    & \multicolumn{4}{p{0.6\linewidth}}{
        \cellcolor[rgb]{0.9,0.9,0.9}{
            \makecell[{{p{\linewidth}}}]{
                \texttt{\tiny{[P1$\rangle$GM]}}
                \texttt{```} \\
\texttt{MUOVERE: O, (4, 2)} \\
\texttt{```} \\
            }
        }
    }
    & & \\ \\

    \theutterance \stepcounter{utterance}  
    & & & \multicolumn{2}{p{0.3\linewidth}}{
        \cellcolor[rgb]{0.9,0.9,0.9}{
            \makecell[{{p{\linewidth}}}]{
                \texttt{\tiny{[GM$|$GM]}}
                \texttt{```} \\
\texttt{MUOVERE: O, (4, 2)} \\
\texttt{```} \\
            }
        }
    }
    & & \\ \\

    \theutterance \stepcounter{utterance}  
    & & & \multicolumn{2}{p{0.3\linewidth}}{
        \cellcolor[rgb]{0.9,0.9,0.9}{
            \makecell[{{p{\linewidth}}}]{
                \texttt{\tiny{[GM$|$GM]}}
                \texttt{O moved to (4, 2)} \\
            }
        }
    }
    & & \\ \\

    \theutterance \stepcounter{utterance}  
    & & \multicolumn{4}{p{0.6\linewidth}}{
        \cellcolor[rgb]{0.9,0.9,0.9}{
            \makecell[{{p{\linewidth}}}]{
                \texttt{\tiny{[GM$\rangle$P2]}}
                \texttt{Il tuo messaggio è stato trasmesso all'altro giocatore.} \\
\\ 
\\ 
\texttt{Stai giocando il round 8 di un massimo di 28.} \\
\texttt{Avete accumulato collettivamente 13 di 16 penalità.} \\
\texttt{L'altro giocatore ha spostato un oggetto sulla sua griglia.} \\
\\ 
\\ 
\texttt{Qual è il tuo prossimo comando?} \\
            }
        }
    }
    & \\ \\

    \theutterance \stepcounter{utterance}  
    & & & \multicolumn{4}{p{0.6\linewidth}}{
        \cellcolor[rgb]{0.9,0.9,0.9}{
            \makecell[{{p{\linewidth}}}]{
                \texttt{\tiny{[GM$\langle$P2]}}
                \texttt{Ora sposto **O** da (4,1) a (4,2) per allinearlo all'obiettivo:} \\
\\ 
\texttt{```} \\
\texttt{MUOVERE: O, (4, 2)} \\
\texttt{```} \\
            }
        }
    }
     \\ \\

    \theutterance \stepcounter{utterance}  
    & & & \multicolumn{2}{p{0.3\linewidth}}{
        \cellcolor[rgb]{0.9,0.9,0.9}{
            \makecell[{{p{\linewidth}}}]{
                \texttt{\tiny{[GM$|$GM]}}
                \texttt{Ora sposto **O** da (4,1) a (4,2) per allinearlo all'obiettivo:} \\
\\ 
\texttt{```} \\
\texttt{MUOVERE: O, (4, 2)} \\
\texttt{```} \\
            }
        }
    }
    & & \\ \\

    \theutterance \stepcounter{utterance}  
    & & & \multicolumn{2}{p{0.3\linewidth}}{
        \cellcolor[rgb]{0.9,0.9,0.9}{
            \makecell[{{p{\linewidth}}}]{
                \texttt{\tiny{[GM$|$GM]}}
                \texttt{Invalid format: head is not empty: 'Ora sposto **O** da (4,1) a (4,2) per allinearlo all'obiettivo:} \\
\\ 
\\ 
\texttt{'} \\
            }
        }
    }
    & & \\ \\

    \theutterance \stepcounter{utterance}  
    & & \multicolumn{4}{p{0.6\linewidth}}{
        \cellcolor[rgb]{0.9,0.9,0.9}{
            \makecell[{{p{\linewidth}}}]{
                \texttt{\tiny{[GM$\rangle$P2]}}
                \texttt{Penalità: Il tuo messaggio non deve contenere nulla prima del comando!} \\
\\ 
\texttt{Assicurati che la tua risposta contenga solo `DIRE: <MESSAGGIO>` o `MUOVERE: <OGGETTO>, (<X>, <Y>)`, e nient'altro!} \\
\\ 
\texttt{Avete accumulato collettivamente 14 di 16 penalità.} \\
\texttt{Per favore riprova!} \\
            }
        }
    }
    & \\ \\

    \theutterance \stepcounter{utterance}  
    & & & \multicolumn{4}{p{0.6\linewidth}}{
        \cellcolor[rgb]{0.9,0.9,0.9}{
            \makecell[{{p{\linewidth}}}]{
                \texttt{\tiny{[GM$\langle$P2]}}
                \texttt{```} \\
\texttt{MUOVERE: O, (4, 2)} \\
\texttt{```} \\
            }
        }
    }
     \\ \\

    \theutterance \stepcounter{utterance}  
    & & & \multicolumn{2}{p{0.3\linewidth}}{
        \cellcolor[rgb]{0.9,0.9,0.9}{
            \makecell[{{p{\linewidth}}}]{
                \texttt{\tiny{[GM$|$GM]}}
                \texttt{```} \\
\texttt{MUOVERE: O, (4, 2)} \\
\texttt{```} \\
            }
        }
    }
    & & \\ \\

    \theutterance \stepcounter{utterance}  
    & & & \multicolumn{2}{p{0.3\linewidth}}{
        \cellcolor[rgb]{0.9,0.9,0.9}{
            \makecell[{{p{\linewidth}}}]{
                \texttt{\tiny{[GM$|$GM]}}
                \texttt{O moved to (4, 2)} \\
            }
        }
    }
    & & \\ \\

    \theutterance \stepcounter{utterance}  
    & & \multicolumn{4}{p{0.6\linewidth}}{
        \cellcolor[rgb]{0.9,0.9,0.9}{
            \makecell[{{p{\linewidth}}}]{
                \texttt{\tiny{[P1$\langle$GM]}}
                \texttt{Spostato 'O' in (4,2) con successo. La tua griglia aggiornata appare così:} \\
\texttt{```} \\
\texttt{ 1234567} \\
\texttt{╔═╤══╤══╗} \\
\texttt{║◌│◌◌│◌◌║ 1} \\
\texttt{╟─┤◌O│◌S║ 2} \\
\texttt{║◌│TD├──╢ 3} \\
\texttt{╟─┘◌◌│◌◌║ 4} \\
\texttt{║◌◌◌P│◌A║ 5} \\
\texttt{║M◌┌─┼──╢ 6} \\
\texttt{║◌◌│◌│◌◌║ 7} \\
\texttt{╚══╧═╧══╝} \\
\texttt{```} \\
\texttt{Stai giocando il round 8 di un massimo di 28.} \\
\texttt{Avete accumulato collettivamente 14 di 16 penalità.} \\
\texttt{L'altro giocatore ha spostato un oggetto sulla sua griglia.} \\
\\ 
\\ 
\texttt{Qual è il tuo prossimo comando?} \\
            }
        }
    }
    & \\ \\

    \theutterance \stepcounter{utterance}  
    & \multicolumn{4}{p{0.6\linewidth}}{
        \cellcolor[rgb]{0.9,0.9,0.9}{
            \makecell[{{p{\linewidth}}}]{
                \texttt{\tiny{[P1$\rangle$GM]}}
                \texttt{```} \\
\texttt{DIRE: O è ora in (4,2). Prossimo obiettivo: T da (3,3) a (6,2). Prima devo liberare (6,2), attualmente occupata da S. Suggerisci dove spostare S?} \\
\texttt{```} \\
            }
        }
    }
    & & \\ \\

    \theutterance \stepcounter{utterance}  
    & & & \multicolumn{2}{p{0.3\linewidth}}{
        \cellcolor[rgb]{0.9,0.9,0.9}{
            \makecell[{{p{\linewidth}}}]{
                \texttt{\tiny{[GM$|$GM]}}
                \texttt{```} \\
\texttt{DIRE: O è ora in (4,2). Prossimo obiettivo: T da (3,3) a (6,2). Prima devo liberare (6,2), attualmente occupata da S. Suggerisci dove spostare S?} \\
\texttt{```} \\
            }
        }
    }
    & & \\ \\

    \theutterance \stepcounter{utterance}  
    & & \multicolumn{4}{p{0.6\linewidth}}{
        \cellcolor[rgb]{0.9,0.9,0.9}{
            \makecell[{{p{\linewidth}}}]{
                \texttt{\tiny{[GM$\rangle$P2]}}
                \texttt{Spostato 'O' in (4,2) con successo. La tua griglia aggiornata appare così:} \\
\texttt{```} \\
\texttt{ 1234567} \\
\texttt{╔═╤══╤══╗} \\
\texttt{║◌│◌◌│T◌║ 1} \\
\texttt{╟─┤◌O│◌◌║ 2} \\
\texttt{║◌│◌◌├──╢ 3} \\
\texttt{╟─┘◌D│◌◌║ 4} \\
\texttt{║P◌A◌│◌S║ 5} \\
\texttt{║◌◌┌─┼──╢ 6} \\
\texttt{║M◌│◌│◌◌║ 7} \\
\texttt{╚══╧═╧══╝} \\
\texttt{```} \\
\texttt{Stai giocando il round 9 di un massimo di 28.} \\
\texttt{Avete accumulato collettivamente 14 di 16 penalità.} \\
\texttt{L'altro giocatore ha inviato questo messaggio:} \\
\\ 
\\ 
\texttt{"O è ora in (4,2). Prossimo obiettivo: T da (3,3) a (6,2). Prima devo liberare (6,2), attualmente occupata da S. Suggerisci dove spostare S?"} \\
\\ 
\\ 
\texttt{Qual è il tuo prossimo comando?} \\
            }
        }
    }
    & \\ \\

    \theutterance \stepcounter{utterance}  
    & & & \multicolumn{4}{p{0.6\linewidth}}{
        \cellcolor[rgb]{0.9,0.9,0.9}{
            \makecell[{{p{\linewidth}}}]{
                \texttt{\tiny{[GM$\langle$P2]}}
                \texttt{```} \\
\texttt{DIRE: Sul mio tabellone S è in (7,5). Puoi spostarlo in (7,2) o (7,3) se libere. Io sposto A da (3,5) a (5,5) per liberare (3,5) per T.} \\
\texttt{```} \\
            }
        }
    }
     \\ \\

    \theutterance \stepcounter{utterance}  
    & & & \multicolumn{2}{p{0.3\linewidth}}{
        \cellcolor[rgb]{0.9,0.9,0.9}{
            \makecell[{{p{\linewidth}}}]{
                \texttt{\tiny{[GM$|$GM]}}
                \texttt{```} \\
\texttt{DIRE: Sul mio tabellone S è in (7,5). Puoi spostarlo in (7,2) o (7,3) se libere. Io sposto A da (3,5) a (5,5) per liberare (3,5) per T.} \\
\texttt{```} \\
            }
        }
    }
    & & \\ \\

    \theutterance \stepcounter{utterance}  
    & & \multicolumn{4}{p{0.6\linewidth}}{
        \cellcolor[rgb]{0.9,0.9,0.9}{
            \makecell[{{p{\linewidth}}}]{
                \texttt{\tiny{[P1$\langle$GM]}}
                \texttt{Il tuo messaggio è stato trasmesso all'altro giocatore.} \\
\\ 
\\ 
\texttt{Stai giocando il round 9 di un massimo di 28.} \\
\texttt{Avete accumulato collettivamente 14 di 16 penalità.} \\
\texttt{L'altro giocatore ha inviato questo messaggio:} \\
\\ 
\\ 
\texttt{"Sul mio tabellone S è in (7,5). Puoi spostarlo in (7,2) o (7,3) se libere. Io sposto A da (3,5) a (5,5) per liberare (3,5) per T."} \\
\\ 
\\ 
\texttt{Qual è il tuo prossimo comando?} \\
            }
        }
    }
    & \\ \\

    \theutterance \stepcounter{utterance}  
    & \multicolumn{4}{p{0.6\linewidth}}{
        \cellcolor[rgb]{0.9,0.9,0.9}{
            \makecell[{{p{\linewidth}}}]{
                \texttt{\tiny{[P1$\rangle$GM]}}
                \texttt{```} \\
\texttt{MUOVERE: S, (7, 2)} \\
\texttt{```} \\
            }
        }
    }
    & & \\ \\

    \theutterance \stepcounter{utterance}  
    & & & \multicolumn{2}{p{0.3\linewidth}}{
        \cellcolor[rgb]{0.9,0.9,0.9}{
            \makecell[{{p{\linewidth}}}]{
                \texttt{\tiny{[GM$|$GM]}}
                \texttt{```} \\
\texttt{MUOVERE: S, (7, 2)} \\
\texttt{```} \\
            }
        }
    }
    & & \\ \\

    \theutterance \stepcounter{utterance}  
    & & & \multicolumn{2}{p{0.3\linewidth}}{
        \cellcolor[rgb]{0.9,0.9,0.9}{
            \makecell[{{p{\linewidth}}}]{
                \texttt{\tiny{[GM$|$GM]}}
                \texttt{Penalità: (7,2) non è vuoto, ma contiene 'S'.} \\
\\ 
\texttt{Avete accumulato collettivamente 15 di 16 penalità.} \\
\texttt{Per favore riprova!} \\
            }
        }
    }
    & & \\ \\

    \theutterance \stepcounter{utterance}  
    & & \multicolumn{4}{p{0.6\linewidth}}{
        \cellcolor[rgb]{0.9,0.9,0.9}{
            \makecell[{{p{\linewidth}}}]{
                \texttt{\tiny{[P1$\langle$GM]}}
                \texttt{Penalità: (7,2) non è vuoto, ma contiene 'S'.} \\
\\ 
\texttt{Avete accumulato collettivamente 15 di 16 penalità.} \\
\texttt{Per favore riprova!} \\
            }
        }
    }
    & \\ \\

    \theutterance \stepcounter{utterance}  
    & \multicolumn{4}{p{0.6\linewidth}}{
        \cellcolor[rgb]{0.9,0.9,0.9}{
            \makecell[{{p{\linewidth}}}]{
                \texttt{\tiny{[P1$\rangle$GM]}}
                \texttt{```} \\
\texttt{DIRE: (7,2) è già occupata da S. Provo a spostare S in (7,3) se è libera. Confermi?} \\
\texttt{```} \\
            }
        }
    }
    & & \\ \\

    \theutterance \stepcounter{utterance}  
    & & & \multicolumn{2}{p{0.3\linewidth}}{
        \cellcolor[rgb]{0.9,0.9,0.9}{
            \makecell[{{p{\linewidth}}}]{
                \texttt{\tiny{[GM$|$GM]}}
                \texttt{```} \\
\texttt{DIRE: (7,2) è già occupata da S. Provo a spostare S in (7,3) se è libera. Confermi?} \\
\texttt{```} \\
            }
        }
    }
    & & \\ \\

    \theutterance \stepcounter{utterance}  
    & & \multicolumn{4}{p{0.6\linewidth}}{
        \cellcolor[rgb]{0.9,0.9,0.9}{
            \makecell[{{p{\linewidth}}}]{
                \texttt{\tiny{[GM$\rangle$P2]}}
                \texttt{Il tuo messaggio è stato trasmesso all'altro giocatore.} \\
\\ 
\\ 
\texttt{Stai giocando il round 10 di un massimo di 28.} \\
\texttt{Avete accumulato collettivamente 15 di 16 penalità.} \\
\texttt{L'altro giocatore ha inviato questo messaggio:} \\
\\ 
\\ 
\texttt{"(7,2) è già occupata da S. Provo a spostare S in (7,3) se è libera. Confermi?"} \\
\\ 
\\ 
\texttt{Qual è il tuo prossimo comando?} \\
            }
        }
    }
    & \\ \\

    \theutterance \stepcounter{utterance}  
    & & & \multicolumn{4}{p{0.6\linewidth}}{
        \cellcolor[rgb]{0.9,0.9,0.9}{
            \makecell[{{p{\linewidth}}}]{
                \texttt{\tiny{[GM$\langle$P2]}}
                \texttt{```} \\
\texttt{MUOVERE: A, (5, 5)} \\
\texttt{```} \\
            }
        }
    }
     \\ \\

    \theutterance \stepcounter{utterance}  
    & & & \multicolumn{2}{p{0.3\linewidth}}{
        \cellcolor[rgb]{0.9,0.9,0.9}{
            \makecell[{{p{\linewidth}}}]{
                \texttt{\tiny{[GM$|$GM]}}
                \texttt{```} \\
\texttt{MUOVERE: A, (5, 5)} \\
\texttt{```} \\
            }
        }
    }
    & & \\ \\

    \theutterance \stepcounter{utterance}  
    & & & \multicolumn{2}{p{0.3\linewidth}}{
        \cellcolor[rgb]{0.9,0.9,0.9}{
            \makecell[{{p{\linewidth}}}]{
                \texttt{\tiny{[GM$|$GM]}}
                \texttt{Penalità: (5,5) non è vuoto, ma contiene '│'.} \\
\\ 
\texttt{Avete accumulato collettivamente 16 di 16 penalità.} \\
\texttt{Per favore riprova!} \\
            }
        }
    }
    & & \\ \\

    \theutterance \stepcounter{utterance}  
    & & \multicolumn{4}{p{0.6\linewidth}}{
        \cellcolor[rgb]{0.9,0.9,0.9}{
            \makecell[{{p{\linewidth}}}]{
                \texttt{\tiny{[GM$\rangle$P2]}}
                \texttt{Penalità: (5,5) non è vuoto, ma contiene '│'.} \\
\\ 
\texttt{Avete accumulato collettivamente 16 di 16 penalità.} \\
\texttt{Per favore riprova!} \\
            }
        }
    }
    & \\ \\

    \theutterance \stepcounter{utterance}  
    & & & \multicolumn{4}{p{0.6\linewidth}}{
        \cellcolor[rgb]{0.9,0.9,0.9}{
            \makecell[{{p{\linewidth}}}]{
                \texttt{\tiny{[GM$\langle$P2]}}
                \texttt{**Partita persa!** Abbiamo raggiunto il limite massimo di 16 penalità collettive.} \\
\\ 
\texttt{Per evitare ulteriori errori, propongo di terminare qui e analizzare gli sbagli:} \\
\texttt{1. **(5,5)** contiene un bordo ('│'), non era una cella valida.} \\
\texttt{2. Dovevamo spostare **A** in **(6,5)** (vuota nel mio tabellone).} \\
\\ 
\texttt{**Comando corretto per riprovare (se la partita continuasse):**} \\
\texttt{```} \\
\texttt{MUOVERE: A, (6, 5)} \\
\texttt{```} \\
\\ 
\texttt{Grazie per la partita! Vuoi ricominciare da capo o chiudere qui?} \\
            }
        }
    }
     \\ \\

    \theutterance \stepcounter{utterance}  
    & & & \multicolumn{2}{p{0.3\linewidth}}{
        \cellcolor[rgb]{0.9,0.9,0.9}{
            \makecell[{{p{\linewidth}}}]{
                \texttt{\tiny{[GM$|$GM]}}
                \texttt{**Partita persa!** Abbiamo raggiunto il limite massimo di 16 penalità collettive.} \\
\\ 
\texttt{Per evitare ulteriori errori, propongo di terminare qui e analizzare gli sbagli:} \\
\texttt{1. **(5,5)** contiene un bordo ('│'), non era una cella valida.} \\
\texttt{2. Dovevamo spostare **A** in **(6,5)** (vuota nel mio tabellone).} \\
\\ 
\texttt{**Comando corretto per riprovare (se la partita continuasse):**} \\
\texttt{```} \\
\texttt{MUOVERE: A, (6, 5)} \\
\texttt{```} \\
\\ 
\texttt{Grazie per la partita! Vuoi ricominciare da capo o chiudere qui?} \\
            }
        }
    }
    & & \\ \\

    \theutterance \stepcounter{utterance}  
    & & & \multicolumn{2}{p{0.3\linewidth}}{
        \cellcolor[rgb]{0.9,0.9,0.9}{
            \makecell[{{p{\linewidth}}}]{
                \texttt{\tiny{[GM$|$GM]}}
                \texttt{Invalid format: head and tail are not empty} \\
\texttt{head: '**Partita persa!** Abbiamo raggiunto il limite massimo di 16 penalità collettive.} \\
\\ 
\texttt{Per evitare ulteriori errori, propongo di terminare qui e analizzare gli sbagli:} \\
\texttt{1. **(5,5)** contiene un bordo ('│'), non era una cella valida.} \\
\texttt{2. Dovevamo spostare **A** in **(6,5)** (vuota nel mio tabellone).} \\
\\ 
\texttt{**Comando corretto per riprovare (se la partita continuasse):**} \\
\\ 
\texttt{'} \\
\texttt{tail: '} \\
\\ 
\\ 
\texttt{Grazie per la partita! Vuoi ricominciare da capo o chiudere qui?'} \\
            }
        }
    }
    & & \\ \\

    \theutterance \stepcounter{utterance}  
    & & & \multicolumn{2}{p{0.3\linewidth}}{
        \cellcolor[rgb]{0.9,0.9,0.9}{
            \makecell[{{p{\linewidth}}}]{
                \texttt{\tiny{[GM$|$GM]}}
                \texttt{Maximum number of penalties exceeded} \\
            }
        }
    }
    & & \\ \\

    \theutterance \stepcounter{utterance}  
    & & & \multicolumn{2}{p{0.3\linewidth}}{
        \cellcolor[rgb]{0.9,0.9,0.9}{
            \makecell[{{p{\linewidth}}}]{
                \texttt{\tiny{[GM$|$GM]}}
                \texttt{Initial states:} \\
\texttt{```} \\
\texttt{Player 1:} \\
\texttt{ 1234567} \\
\texttt{╔═╤══╤══╗} \\
\texttt{║◌│◌◌│◌◌║ 1} \\
\texttt{╟─┤◌◌│◌S║ 2} \\
\texttt{║◌│TD├──╢ 3} \\
\texttt{╟─┘◌◌│O◌║ 4} \\
\texttt{║◌◌PM│◌A║ 5} \\
\texttt{║◌◌┌─┼──╢ 6} \\
\texttt{║◌◌│◌│◌◌║ 7} \\
\texttt{╚══╧═╧══╝} \\
\\ 
\\ 
\texttt{Player 2:} \\
\texttt{ 1234567} \\
\texttt{╔═╤══╤══╗} \\
\texttt{║◌│◌O│T◌║ 1} \\
\texttt{╟─┤◌◌│◌◌║ 2} \\
\texttt{║◌│◌◌├──╢ 3} \\
\texttt{╟─┘◌D│◌◌║ 4} \\
\texttt{║◌SA◌│◌◌║ 5} \\
\texttt{║PM┌─┼──╢ 6} \\
\texttt{║◌◌│◌│◌◌║ 7} \\
\texttt{╚══╧═╧══╝} \\
\\ 
\texttt{```} \\
            }
        }
    }
    & & \\ \\

    \theutterance \stepcounter{utterance}  
    & & & \multicolumn{2}{p{0.3\linewidth}}{
        \cellcolor[rgb]{0.9,0.9,0.9}{
            \makecell[{{p{\linewidth}}}]{
                \texttt{\tiny{[GM$|$GM]}}
                \texttt{End states:} \\
\texttt{```} \\
\texttt{Player 1:} \\
\texttt{ 1234567} \\
\texttt{╔═╤══╤══╗} \\
\texttt{║◌│◌◌│◌◌║ 1} \\
\texttt{╟─┤◌O│◌S║ 2} \\
\texttt{║◌│TD├──╢ 3} \\
\texttt{╟─┘◌◌│◌◌║ 4} \\
\texttt{║◌◌◌P│◌A║ 5} \\
\texttt{║M◌┌─┼──╢ 6} \\
\texttt{║◌◌│◌│◌◌║ 7} \\
\texttt{╚══╧═╧══╝} \\
\\ 
\\ 
\texttt{Player 2:} \\
\texttt{ 1234567} \\
\texttt{╔═╤══╤══╗} \\
\texttt{║◌│◌◌│T◌║ 1} \\
\texttt{╟─┤◌O│◌◌║ 2} \\
\texttt{║◌│◌◌├──╢ 3} \\
\texttt{╟─┘◌D│◌◌║ 4} \\
\texttt{║P◌A◌│◌S║ 5} \\
\texttt{║◌◌┌─┼──╢ 6} \\
\texttt{║M◌│◌│◌◌║ 7} \\
\texttt{╚══╧═╧══╝} \\
\\ 
\texttt{```} \\
            }
        }
    }
    & & \\ \\

    \theutterance \stepcounter{utterance}  
    & & & \multicolumn{2}{p{0.3\linewidth}}{
        \cellcolor[rgb]{0.9,0.9,0.9}{
            \makecell[{{p{\linewidth}}}]{
                \texttt{\tiny{[GM$|$GM]}}
                \texttt{* success: False} \\
\texttt{* lose: False} \\
\texttt{* aborted: True} \\
\texttt{{-}{-}{-}{-}{-}{-}{-}} \\
\texttt{* Shifts: 4.00} \\
\texttt{* Max Shifts: 12.00} \\
\texttt{* Min Shifts: 6.00} \\
\texttt{* End Distance Sum: 15.61} \\
\texttt{* Init Distance Sum: 22.51} \\
\texttt{* Expected Distance Sum: 29.33} \\
\texttt{* Penalties: 17.00} \\
\texttt{* Max Penalties: 16.00} \\
\texttt{* Rounds: 9.00} \\
\texttt{* Max Rounds: 28.00} \\
\texttt{* Object Count: 7.00} \\
            }
        }
    }
    & & \\ \\

\end{supertabular}
}

\end{document}
